% !TeX root = ../../Book1.tex
% 纲要标题：## 第21章　分圆扩张、根式与尺规作图
% 纲要位置：计划纲要.md，第 736 行。

\chapter[分圆扩张、根式与尺规作图]{分圆扩张、\mbox{根式与尺规作图}}
\label{b1:ch:21}
\BookChapterNumber{b1:ch:21}

\begin{chapterabstract}
单位根把正多边形作图转成分圆域的结构问题；根式求解则与多项式 Galois 群的可解性相连。本章建立分圆域及根式可解性理论，比较低次方程与一般五次方程，并以二次扩张塔判断尺规可作性。
\end{chapterabstract}

\begin{learninggoals}
\begin{itemize}
\item 从全部单位根中分离本原单位根，证明并运用分圆多项式的不可约性。
\item 区分根式塔、分裂域和正规闭包，掌握根式可解性判据的两个方向。
\item 解释一般方程的独立系数含义，避免将一般五次结论误用于所有特殊五次方程。
\item 把尺规操作与实二次扩张塔互相转换，判断典型长度、角和正多边形是否可作图。
\end{itemize}
\end{learninggoals}

正五边形与正七边形都可以由单位圆上的根来描述。前者涉及的分圆域在有理数域上次数为 \(4\)，且其 Galois 群为循环群 \(C_4\)，由指数为二的子群可得到两次二次扩张；后者的相应次数为 \(6\)，不可能包含在二次扩张塔的末端。这解释了为什么尺规能作正五边形，却不能作正七边形。这里要先弄清单位根生成了什么域、域的自同构怎样排列这些根，才能把几何作图翻译成域扩张问题。类似地，讨论方程能否用根式求解，也须把允许的求根步骤同相应 Galois 群的结构联系起来。

\ProseParagraphBreak

本章使用\cref{b1:thm:finite-galois-correspondence}、\cref{b1:thm:kummer-cyclic}和可解群的性质；分圆不可约性还需\cref{b1:thm:gauss-descent}。一般五次方程的结论须另算独立参数多项式的 Galois 群，具体多项式 \(X^5-X-1\) 则使用\cref{b1:lem:reduction-cycle-type}和置换群论证。作图判据以平面几何、二次方程和二次扩张塔为基础；正多边形还需分圆域与 Galois 对应。圆周率超越性的证明概要另用复指数函数、路径积分与代数整数范数。

\ProseParagraphBreak
单位根与根式可解性可读 Roman 的《Field Theory》\cite[Sections 11.1--11.2, 13.2--13.7]{Roman2006FieldTheory}。三次方程、实根和对称多项式的中文基础材料可见席南华的《基础代数》第一卷\cite[第7章]{Xi2016BasicAlgebraI}。尺规作图末尾的超越性概要依据 Lindemann 的原文\cite[\S\S 1--4, pp. 214--225]{Lindemann1882Pi}，那里使用的积分估计与本章的代数作图判据分别承担不同的论证。

% !TeX root = ../../Book1.tex
% 纲要标题：### 21.1 分圆多项式
% 纲要位置：计划纲要.md，第 738 行。

\section{分圆多项式}
\label{b1:sec:2101}

上一章在“基域已经含有所需单位根”的条件下讨论开方。现在把基域取为 \(\mathbb Q\)，研究加入单位根本身造成的扩张。需要区分两件事：\(X^n-1\) 收集全部 \(n\) 次单位根，而分圆多项式只收集乘法阶恰为 \(n\) 的那些根。

% 纲要内容（仅作写作计划，尚非教材正文）：
%
% * 单位根与分圆多项式
% * 分圆多项式的不可约性
% * 分圆域及其 Galois 群；一般特征零基域上的单位根扩张只嵌入相应单位群
%

\subsection{单位根与分圆多项式}
\label{b1:subsec:2101:01}

取 \(\zeta_n=e^{2\pi i/n}\)。它是一个本原 \(n\) 次单位根；所有本原 \(n\) 次单位根为 \(\zeta_n^a\)，其中 \(1\leq a\leq n\)、\(\gcd(a,n)=1\)。

\begin{definition}\label{b1:def:cyclotomic-polynomial}
设 \(n\geq1\) 为整数，并在 \(\mathbb C\) 中选取一个本原 \(n\) 次单位根 \(\zeta_n\)。多项式
\[
\Phi_n(X)\coloneqq\prod_{\substack{1\leq a\leq n\\\gcd(a,n)=1}}(X-\zeta_n^a)
\]
不依赖本原单位根 \(\zeta_n\) 的选择，称为第 \(n\) 个\term[fenyuan-duoxiangshi]{分圆多项式}[cyclotomic polynomial]。其次数为 Euler 函数 \(\varphi(n)\)，即模 \(n\) 可逆剩余类的个数。
\end{definition}
\symidx{\(\Phi_n(X)\)：第 \(n\) 个分圆多项式}

若改取另一个本原单位根 \(\zeta_n'=\zeta_n^b\)，则 \(\gcd(b,n)=1\)。在单位群 \((\mathbb Z/n\mathbb Z)^\times\) 中，乘以 \(b\) 是一个置换，所以 \((\zeta_n')^a=\zeta_n^{ba}\) 仍遍历同一组本原根。这就证明了定义中的乘积不依赖 \(\zeta_n\) 的选择。

\ProseParagraphBreak
在特征零中，\(X^n-1\) 与导数 \(nX^{n-1}\) 互素，所以它的 \(n\) 个根彼此不同。每个根的乘法阶又是 \(n\) 的唯一一个正因子，按阶分组便把全部根不重不漏地分到各个 \(\Phi_d\) 中。因此
\begin{equation}\label{b1:eq:cyclotomic-factorization}
X^n-1=\prod_{d\mid n}\Phi_d(X).
\end{equation}
这还给出 \(\Phi_n\in\mathbb Z[X]\)。对 \(n\) 归纳，\(\Phi_1=X-1\)；若真因子对应的多项式均已知为首一整系数多项式，则用它们的乘积去除 \(X^n-1\)。除式首一，每次消去最高项时只须减去除式的一个整系数单项式倍数，无须除以非单位的首项系数，所以商、余式仍在 \(\mathbb Z[X]\) 中。而\eqref{b1:eq:cyclotomic-factorization}说明在 \(\mathbb C[X]\) 中余式为零、商为 \(\Phi_n\)；除法的唯一性便给出结论。

\ProseParagraphBreak
例如 \(\Phi_2=X+1\)、\(\Phi_3=X^2+X+1\)、\(\Phi_4=X^2+1\)、\(\Phi_6=X^2-X+1\)。对素数 \(p\)，有 \(\Phi_p=1+X+\cdots+X^{p-1}\)。不过，整系数性本身不能说明不可约；下一小节需要另外论证所有本原根互为有理数域上的共轭。

\subsection{分圆多项式的不可约性}
\label{b1:subsec:2101:02}

要证明 \(\Phi_n\) 不可约，只须证明一个本原根的最小多项式已经包含全部本原根。我们将证明，它的根取任意不整除 \(n\) 的素数次幂后，仍是同一个最小多项式的根；将互素指数分解成素因子，就能从 \(\zeta_n\) 逐次到达每个 \(\zeta_n^a\)。

\begin{theorem}\label{b1:thm:cyclotomic-irreducible}
设 \(n\geq1\) 为整数，\(\zeta_n\in\mathbb C\) 为本原 \(n\) 次单位根。则第 \(n\) 个分圆多项式 \(\Phi_n(X)\) 在 \(\mathbb Q[X]\) 中不可约。因此 \(\bigl[\mathbb Q(\zeta_n):\mathbb Q\bigr]=\varphi(n)\)。
\end{theorem}
\begin{proof}
设 \(f\) 为 \(\zeta_n\) 在 \(\mathbb Q\) 上的首一最小多项式。它是首一整系数多项式 \(\Phi_n\) 的因子，由\cref{b1:thm:gauss-descent}，可写 \(\Phi_n=fg\)，其中 \(f,g\in\mathbb Z[X]\) 都首一。

设 \(\alpha\) 是 \(f\) 的根，\(q\) 为不整除 \(n\) 的素数。由 \(f\mid\Phi_n\)，\(\alpha\) 的乘法阶为 \(n\)。又因 \(\gcd(q,n)=1\)，\cref{b1:prop:power-order}给出
\[
\operatorname{ord}(\alpha^q)=\frac{n}{\gcd(n,q)}=n,
\]
故 \(\alpha^q\) 仍是 \(\Phi_n\) 的根。若它不是 \(f\) 的根，由 \(\Phi_n=fg\) 就有 \(g(\alpha^q)=0\)。\(f\) 不可约且零化 \(\alpha\)，所以它也是 \(\alpha\) 的首一最小多项式，从而 \(f(X)\mid g(X^q)\) 于 \(\mathbb Q[X]\)。首一除法又说明这个整除关系的商在 \(\mathbb Z[X]\) 中。

把这个整系数整除关系模 \(q\) 约化，并以横线表示所得多项式。特征 \(q\) 的二项式公式以及 \(c^q=c\) 对 \(c\in\mathbb F_q\) 的成立给出
\[
\bar f(X)\mid\bar g(X^q)=\bar g(X)^q.
\]
\(f\) 首一，故约化不会降低次数，\(\bar f\) 仍有正次数，可取它的一个不可约因子 \(h\)。域上的一元多项式环 \(\mathbb F_q[X]\) 是唯一分解整环，\cref{b1:prop:ufd-prime}说明其中的不可约元是素元。因此 \(h\mid\bar g^q\) 迫使 \(h\mid\bar g\)。于是 \(h\) 同时整除 \(\bar f\) 和 \(\bar g\)，给出 \(h^2\mid\bar f\bar g=\overline{\Phi_n}\)。再把\eqref{b1:eq:cyclotomic-factorization}模 \(q\) 约化，得到 \(h^2\mid X^n-1\)。但 \(q\nmid n\)，所以导数 \(nX^{n-1}\) 的系数 \(n\) 非零，而 \(X\) 不整除 \(X^n-1\)；两多项式互素，不可能有重不可约因子，矛盾。这证明了 \(\alpha^q\) 仍为 \(f\) 的根。

任意与 \(n\) 互素的正整数 \(a>1\) 可写成 \(a=q_1\cdots q_s\)，其中每个素数 \(q_j\) 都不整除 \(n\)。从已知的根 \(\zeta_n\) 出发，刚证的结论依次说明
\[
\zeta_n^{q_1},\quad \zeta_n^{q_1q_2},\quad\ldots,\quad\zeta_n^{q_1\cdots q_s}=\zeta_n^a
\]
都是 \(f\) 的根；\(a=1\) 时无需这一步。因此 \(f\) 含有全部 \(\varphi(n)\) 个互异的本原根，而它又整除同为首一的 \(\Phi_n\)，只能有 \(f=\Phi_n\)。次数结论来自\cref{b1:prop:minimal-polynomial-degree}。
\end{proof}

\subsection{分圆域及其 Galois 群}
\label{b1:subsec:2101:03}

\begin{theorem}\label{b1:thm:cyclotomic-galois-group}
设 \(n\geq1\) 为整数，\(\zeta_n\in\mathbb C\) 为本原 \(n\) 次单位根。则\term[fenyuan-yu]{分圆域}[cyclotomic field] \(\mathbb Q(\zeta_n)\) 是 \(X^n-1\) 在 \(\mathbb Q\) 上的分裂域，且
\[
\operatorname{Gal}\bigl(\mathbb Q(\zeta_n)/\mathbb Q\bigr)
\cong(\mathbb Z/n\mathbb Z)^\times,
\qquad \sigma_a(\zeta_n)=\zeta_n^a.
\]
\end{theorem}
\begin{proof}
全部 \(n\) 次单位根都是 \(\zeta_n\) 的幂，所以 \(\mathbb Q(\zeta_n)\) 正是 \(X^n-1\) 的分裂域。它有限且正规，特征零又保证可分，因而为有限 Galois 扩张。自同构保持元素乘法阶，所以每个自同构 \(\sigma\) 都满足 \(\sigma(\zeta_n)=\zeta_n^a\)，其中 \(\gcd(a,n)=1\)。剩余类 \(a\bmod n\) 唯一，而且自同构由生成元 \(\zeta_n\) 的像决定，于是得到单射
\[
\operatorname{Gal}\bigl(\mathbb Q(\zeta_n)/\mathbb Q\bigr)
\longrightarrow(\mathbb Z/n\mathbb Z)^\times.
\]
若 \(\sigma(\zeta_n)=\zeta_n^a\)、\(\tau(\zeta_n)=\zeta_n^b\)，则 \((\sigma\tau)(\zeta_n)=\zeta_n^{ab}\)，故这也是群同态。由\cref{b1:prop:galois-automorphism-count,b1:thm:cyclotomic-irreducible}，源群的阶为
\[
\left|\operatorname{Gal}\bigl(\mathbb Q(\zeta_n)/\mathbb Q\bigr)\right|
=[\mathbb Q(\zeta_n):\mathbb Q]=\varphi(n).
\]
目标单位群的阶按 \(\varphi(n)\) 的定义也为 \(\varphi(n)\)，所以单射必为满射，得到所述同构。
\end{proof}

例如 \(\mathbb Q(\zeta_5)/\mathbb Q\) 的群为 \(C_4\)，而 \(\mathbb Q(\zeta_8)/\mathbb Q\) 的群为 \(C_2\times C_2\)，不能把所有分圆扩张都称为循环扩张。“循环扩张”要求 Galois 群循环；\(\mathbb Q(\zeta_8)\) 虽由一个元素生成，却不满足这个要求。

\ProseParagraphBreak
复共轭 \(c\) 总对应于剩余类 \(-1\)。设 \(n\geq3\)，记 \(t=\zeta_n+\zeta_n^{-1}\)，则 \(c(t)=t\)，所以
\[
\mathbb Q(t)\subseteq\mathbb Q(\zeta_n)^{\langle c\rangle}.
\]
同时 \(\mathbb Q(t)\subseteq\mathbb R\)，而 \(\zeta_n\) 非实，并满足
\[
X^2-(\zeta_n+\zeta_n^{-1})X+1=0;
\]
故 \([\mathbb Q(\zeta_n):\mathbb Q(t)]=2\)。子群 \(\langle c\rangle\) 的阶也为 \(2\)，由\eqref{b1:eq:galois-degree-index}，它的不动域同样满足 \([\mathbb Q(\zeta_n):\mathbb Q(\zeta_n)^{\langle c\rangle}]=2\)。结合上述包含关系和塔式公式，得到这两个中间域相等。

\ProseParagraphBreak
把基域 \(\mathbb Q\) 换成任意特征零域 \(K\) 后，自同构还须固定 \(K\) 中的其他元素，因而不能再保证每个互素幂次都能作为 \(\zeta_n\) 的像。不过，确定生成元像的办法仍给出到单位群的单同态，所以 Galois 群仍然交换。

\begin{proposition}\label{b1:prop:cyclotomic-general-base}
设 \(K\) 为特征零域，\(\Omega\) 为 \(K\) 的代数闭包，\(n\geq1\) 为整数，\(\mu_n\subseteq\Omega\) 为全部 \(n\) 次单位根。选取一个本原 \(n\) 次单位根 \(\zeta_n\in\Omega\)，并令 \(L=K(\mu_n)=K(\zeta_n)\)。则 \(L/K\) 为有限 Galois 扩张，映射
\[
\operatorname{Gal}(L/K)\longrightarrow(\mathbb Z/n\mathbb Z)^\times,
\qquad\sigma\longmapsto a\bmod n\quad
\bigl(\sigma(\zeta_n)=\zeta_n^a\bigr)
\]
是单同态。因此 \(\operatorname{Gal}(L/K)\) 是交换群，且 \([L:K]\) 整除 \(\varphi(n)\)。
\end{proposition}
\begin{proof}
特征零保证 \(X^n-1\) 有 \(n\) 个互异的根，它们组成乘法群 \(\mu_n\)。由\cref{b1:lem:finite-multiplicative-cyclic}，该群循环，存在所述的本原单位根，且全部根都是 \(\zeta_n\) 的幂。因此 \(L=K(\zeta_n)\) 是 \(X^n-1\) 的分裂域，为有限 Galois 扩张。

每个 \(\sigma\in\operatorname{Gal}(L/K)\) 保持 \(\zeta_n\) 的乘法阶，所以 \(\sigma(\zeta_n)=\zeta_n^a\)，其中 \(a\bmod n\) 是一个唯一确定的可逆剩余类。生成元的像决定自同构，故所列映射单射；复合对应指数相乘，故它是同态。其像是交换群 \((\mathbb Z/n\mathbb Z)^\times\) 的子群，因而交换，且阶整除 \(\varphi(n)\)。再由有限 Galois 扩张的群阶等于次数，得到次数的整除结论。
\end{proof}

这个单射在一般基域上不必满射。例如 \(K\) 本来就包含 \(\mu_n\) 时，\(L=K\)，Galois 群平凡；\(n\geq3\) 时，单位群却并不平凡。

% !TeX root = ../../Book1.tex
% 纲要标题：### 21.2 根式可解性
% 纲要位置：计划纲要.md，第 744 行。

\section{根式可解性}
\label{b1:sec:2102}

一个根式公式允许有限次四则运算和开方，也允许在计算途中使用复数。域塔把这样的表达式记录为一串有限扩张；Galois 群则检验这一串扩张是否可能存在。本节的等价定理以特征零为主版本，且把“一个根能够表示”与“多项式的全部根能够表示”明确分开。

% 纲要内容（仅作写作计划，尚非教材正文）：
%
% * 根式塔与可解扩张
% * 根式塔的正规可解包络
% * 根式可解性与可解群：集中陈述特征零判据并连续证明两个方向
% * 主证明完成后讨论特征条件与可分性约定
% * 完整的根式可解性等价定理以特征零为主版本：在加入必要单位根后，调用第20.3节的循环扩张根式描述，沿可解群的素数阶循环商逐层构造；正特征情形单独说明与 Artin–Schreier 扩张的差别。
%

\subsection{根式塔与可解扩张}
\label{b1:subsec:2102:01}

\begin{definition}\label{b1:def:radical-extension}
设 \(K\) 为域，\(r\geq0\) 为整数。给定有限域扩张塔
\[
K=K_0\subseteq K_1\subseteq\cdots\subseteq K_r,
\qquad K_i=K_{i-1}(\alpha_i),\quad\alpha_i^{m_i}\in K_{i-1},
\]
其中每个 \(m_i\geq2\) 都是整数。如果各步都满足所列条件，则称这座塔为\term[genshi-ta]{根式塔}[radical tower]，并称末端扩张 \(K_r/K\) 为\term[genshi-kuozhang]{根式扩张}[radical extension]。给定多项式 \(f\in K[X]\)，如果存在一座从 \(K\) 出发的根式塔，其末端域包含 \(f\) 在一个代数闭包中的全部根，则称 \(f\) \term[genshi-kejie]{根式可解}[solvable by radicals]。
\end{definition}

\(\alpha_i=0\) 或本已属于 \(K_{i-1}\) 的步骤可以删去。开方的指数可以是合数，但能够分拆：若 \(\alpha^{rs}=a\in F\)，其中 \(r,s\geq2\)，令 \(\beta=\alpha^r\)，则 \(\beta^s=a\) 且 \(\alpha^r=\beta\)。于是原来一步 \(F\subseteq F(\alpha)\) 可以改成
\[
F\subseteq F(\beta)\subseteq F(\alpha),
\]
分别取 \(s\) 次方根和 \(r\) 次方根。继续分解指数，就得到一条只取素数次方根、末端域相同的塔。这里并不要求每一步正规，也不要求 \(K_r\) 恰好等于分裂域。例如 \(\mathbb Q(\sqrt[3]{2})\) 是根式扩张，但还没有包含 \(X^3-2\) 的全部根；再加入 \(\zeta_3\) 即可。

\begin{definition}\label{b1:def:solvable-extension}
设 \(K\) 的特征为零，\(E/K\) 为有限扩张。若 \(E/K\) 的正规闭包的 Galois 群是可解群，则称 \(E/K\) 为\term[kejie-kuozhang]{可解扩张}[solvable extension]。
\end{definition}

这里的正规闭包可按\cref{b1:prop:normal-closure}取为有限生成元的最小多项式乘积的分裂域，所以仍是有限正规扩张。特征零又保证它可分，因而确为有限 Galois 扩张。

\ProseParagraphBreak
根式扩张要求这个具体的域 \(E\) 本身能作为一座根式塔的末端；可解扩张只要求它的正规闭包具有可解的 Galois 群。若 \(E/K\) 已经是有限 Galois 扩张，后一个定义中的群就是 \(\operatorname{Gal}(E/K)\)。下面的判据将说明，可解性保证 \(E\) 能包含在某座根式塔的末端域中，而这个末端域可以严格大于 \(E\)。稍后的实三次分圆子域正是 \(E/K\) 可解而 \(E/K\) 本身不是根式扩张的例子。

\subsection{根式塔的正规可解包络}
\label{b1:subsec:2102:02}

根式塔的各层未必正规。单位根保证同一个二项式的根只相差一个已知因子，但补齐各层原方程的根还不够：基域上的共轭变换也会移动方程的系数。因此，要把根式塔放进一个 Galois 扩张，在加入每个根式时还须同时补入共轭方程的根。

\ProseParagraphBreak
例如令 \(\alpha=\sqrt{1+\sqrt2}>0\)。根式塔
\[
\mathbb Q\subseteq\mathbb Q(\sqrt2)\subseteq\mathbb Q(\alpha)
\]
的第二步由 \(X^2-(1+\sqrt2)\) 给出；这里 \(\sqrt2=\alpha^2-1\)，所以前一层确实包含在末端域中。这个二项式的两个根 \(\pm\alpha\) 已经都在末端域中，二次单位根也已经属于 \(\mathbb Q\)。可是将 \(\sqrt2\) 送到 \(-\sqrt2\) 的共轭变换，会把方程变成
\[
X^2-(1-\sqrt2)=0.
\]
它的根是非实数，不属于实域 \(\mathbb Q(\alpha)\)。因此补齐原方程的根仍不能得到 \(\mathbb Q\) 上的正规扩张；还须补入共轭方程的根。一般构造中让前一层的自同构遍历，正是为了包含这些随系数变化而出现的方程。

\begin{lemma}\label{b1:lem:radical-normal-envelope}
特征零域 \(K\) 上的任意有限根式塔，都包含在一个有限 Galois 扩张 \(M/K\) 中，使 \(\operatorname{Gal}(M/K)\) 可解。
\end{lemma}
\begin{proof}
若塔中所有域都等于 \(K\)，取 \(M=K\) 即可。以下略去平凡步骤，记 \(K_i=K_{i-1}(\alpha_i)\)、\(\alpha_i^{m_i}=a_i\neq0\)，并固定一个包含末端域的 \(K\) 的代数闭包 \(\Omega\)。我们将在 \(\Omega\) 中逐层构造 \(M_i\)，使 \(K_i\subseteq M_i\)，而 \(M_i/K\) 始终是群可解的有限 Galois 扩张。

取全部 \(m_i\) 的公倍数 \(m\)，先令 \(M_0=K(\mu_m)\)。由\cref{b1:prop:cyclotomic-general-base}，它是 \(X^m-1\) 的有限 Galois 分裂域，其 Galois 群嵌入交换群 \((\mathbb Z/m\mathbb Z)^\times\)。所以这个群本身交换，特别地可解。

归纳设 \(M_{i-1}/K\) 为有限 Galois 扩张，包含 \(K_{i-1}\)，且群可解。令 \(M_i\) 为 \(M_{i-1}\) 上全部多项式
\[
X^{m_i}-\tau(a_i),\qquad\tau\in\operatorname{Gal}(M_{i-1}/K),
\]
在 \(\Omega\) 中的共同分裂域。这族多项式只有有限多个，每一个又只有有限多个代数根，故 \(M_i/M_{i-1}\) 有限，从而 \(M_i/K\) 有限。取 \(\tau=1\) 时的方程，其根 \(\alpha_i\) 必在共同分裂域中，所以 \(K_i\subseteq M_i\)。让 \(\tau\) 遍历的目的，正是把前例中随 \(\sqrt2\mapsto-\sqrt2\) 出现的新方程一并加入。

任取一个 \(K\) 嵌入 \(\iota:M_i\rightarrow\Omega\)。因 \(M_{i-1}/K\) 正规，它在 \(M_{i-1}\) 上的限制为某个 \(\sigma\in\operatorname{Gal}(M_{i-1}/K)\)。若 \(\beta\) 是 \(X^{m_i}-\tau(a_i)\) 的根，则
\[
\iota(\beta)^{m_i}=\iota\bigl(\tau(a_i)\bigr)=(\sigma\tau)(a_i).
\]
因此 \(\iota\) 将这个方程的根集映到指标为 \(\sigma\tau\) 的方程的根集。因 \(\tau(a_i)\neq0\) 且特征为零，两个根集均有 \(m_i\) 个元素，嵌入又是单射，所以该映射是双射；左乘 \(\sigma\) 同时置换全部指标。于是 \(\iota\) 把这些根和 \(M_{i-1}\) 生成的域 \(M_i\) 映到自身且映满。由\cref{b1:thm:normal-embedding-criterion}，\(M_i/K\) 正规；特征零保证可分，故它是有限 Galois 扩张。

对每个所列多项式选一个根 \(\beta_\tau\)。因为 \(m_i\mid m\)，且 \(M_0\subseteq M_{i-1}\)，有 \(\mu_{m_i}\subseteq\mu_m\subseteq M_{i-1}\)。每个方程的全部根因而为 \(\zeta\beta_\tau\)，其中 \(\zeta\in\mu_{m_i}\)，所以
\[
M_i=M_{i-1}\bigl(\beta_\tau\mid\tau\in\operatorname{Gal}(M_{i-1}/K)\bigr).
\]
对 \(\rho\in\operatorname{Gal}(M_i/M_{i-1})\)，因 \(\rho\) 固定 \(\tau(a_i)\)，比值 \(\rho(\beta_\tau)/\beta_\tau\) 属于 \(\mu_{m_i}\)。由此得到映射
\[
\operatorname{Gal}(M_i/M_{i-1})\longrightarrow
\prod_{\tau}\mu_{m_i},\qquad
\rho\longmapsto\bigl(\rho(\beta_\tau)/\beta_\tau\bigr)_\tau
\]
它与\cref{b1:prop:single-radical-cyclic}中的比值映射相同。具体地，\(\rho\) 固定 \(M_{i-1}\)，故固定全部单位根；对另一个相对自同构 \(\eta\)，有
\[
\frac{\rho\eta(\beta_\tau)}{\beta_\tau}
=\rho\Bigl(\frac{\eta(\beta_\tau)}{\beta_\tau}\Bigr)
\frac{\rho(\beta_\tau)}{\beta_\tau}
=\frac{\eta(\beta_\tau)}{\beta_\tau}
\frac{\rho(\beta_\tau)}{\beta_\tau}.
\]
目标群交换，上式便是逐坐标的同态公式。若全部比值为 \(1\)，\(\rho\) 固定全部生成根，也就固定整个 \(M_i\)，故这个同态单射。单位根的直积是交换群，其子群也交换，所以 \(\operatorname{Gal}(M_i/M_{i-1})\) 实际上是交换群。由于 \(M_{i-1}/K\) 已是 Galois 扩张，\cref{b1:thm:galois-quotient}给出限制映射的正合列
\[
1\longrightarrow\operatorname{Gal}(M_i/M_{i-1})
\longrightarrow\operatorname{Gal}(M_i/K)
\longrightarrow\operatorname{Gal}(M_{i-1}/K)\longrightarrow1.
\]
其核是刚证明交换的相对 Galois 群，商是归纳中已知可解的 \(\operatorname{Gal}(M_{i-1}/K)\)。由\cref{b1:thm:solvable-closure}，\(\operatorname{Gal}(M_i/K)\) 可解。这就保持了构造中所需的全部条件；取 \(M=M_r\)，它包含根式塔的末端域，因而得到结论。
\end{proof}

\subsection{根式可解性与可解群}
\label{b1:subsec:2102:04}

\begin{theorem}[根式可解性判据]\label{b1:thm:radical-solvability}
设 \(K\) 的特征为零，\(f\in K[X]\) 非常数，\(L\) 为其分裂域。则 \(f\) 根式可解当且仅当 \(\operatorname{Gal}(L/K)\) 是可解群。
\end{theorem}

由根式塔推出群可解，所用的是正规可解包络及分裂域群作为商群的性质；由群可解构造根式塔，所用的则是素数阶循环因子和 Kummer 定理。后一方向须先加入单位根，才能把循环扩张转成开方。

\begin{proof}
先设 \(f\) 的全部根包含在某条根式塔中。由\cref{b1:lem:radical-normal-envelope}，该塔包含于有限 Galois 扩张 \(M/K\)，其群可解。\(L\) 由 \(f\) 的全部根在 \(K\) 上生成，故 \(L\subseteq M\)。又因 \(L/K\) 是特征零域上多项式的分裂域，它也为有限 Galois 扩张。对 \(M/K\) 的中间域 \(L\) 应用\cref{b1:thm:galois-quotient}，\(\operatorname{Gal}(M/L)\) 正规于 \(\operatorname{Gal}(M/K)\)，限制同态满射到 \(\operatorname{Gal}(L/K)\)，且
\[
\operatorname{Gal}(L/K)\cong
\operatorname{Gal}(M/K)/\operatorname{Gal}(M/L).
\]
由\cref{b1:thm:solvable-closure}，这个商群可解，得到必要性。

反过来，设 \(G=\operatorname{Gal}(L/K)\) 可解。取 \(m=|G|\)。若 \(m=1\)，由 \([L:K]=|G|\) 已有 \(L=K\)，无须开方。以下设 \(m>1\)，选取本原 \(m\) 次单位根 \(\zeta_m\)，令 \(F=K(\mu_m)=K(\zeta_m)\)、\(E=LF\)。因为 \(\zeta_m^m=1\in K\)，这一个特定的 \(m\) 次方根就给出根式扩张 \(F/K\)。\(E/F\) 是 \(f\) 在 \(F\) 上的分裂域，有限且 Galois。对任意 \(\sigma\in\operatorname{Gal}(E/F)\)，因 \(K\subseteq F\)，\(\sigma|_L\) 是一个 \(K\) 嵌入；\(L/K\) 的正规性又保证 \(\sigma(L)=L\)，所以限制给出同态
\[
\operatorname{Gal}(E/F)\hookrightarrow\operatorname{Gal}(L/K),
\]
如果 \(\sigma\) 的限制为恒等，它既固定 \(L\) 又固定 \(F\)，便固定这两个域生成的 \(E=LF\)，所以核平凡。由这个单同态及\cref{b1:thm:solvable-closure}，\(H=\operatorname{Gal}(E/F)\) 可解，且其阶整除 \(m\)。

由\cref{b1:cor:solvable-composition-factors}，有限可解群 \(H\) 有素数阶循环因子的次正规列。其产生过程是：先由\cref{b1:prop:derived-quotient}取交换因子的导出列，再在每个非平凡有限交换因子中取极大真子群，并对阶较小的子群继续做同样的选择。交换群的子群都正规，极大性使所得商为单群；由\cref{b1:prop:abelian-simple}，每个商都是素数阶循环群。通过商映射取这些子群的原像，就把导出列细化为
\[
H=H_0\trianglerighteq H_1\trianglerighteq\cdots\trianglerighteq H_s=1,
\qquad H_{i-1}/H_i\cong C_{\ell_i}.
\]
\cref{b1:thm:finite-galois-correspondence}给出反向域塔 \(F=E_0\subseteq E_1\subseteq\cdots\subseteq E_s=E\)，其中 \(E_i=E^{H_i}\)。相邻的包含 \(H_i\subseteq H_{i-1}\) 因而对应 \(E_{i-1}=E^{H_{i-1}}\subseteq E^{H_i}=E_i\)。为确定这一步的 Galois 群，把基域换成 \(E_{i-1}\)：\(E/E_{i-1}\) 仍是有限 Galois 扩张，且
\[
\operatorname{Gal}(E/E_{i-1})=H_{i-1}.
\]
在这个群中，\(H_i\trianglelefteq H_{i-1}\)，其不动域为 \(E_i\)，所以\cref{b1:thm:galois-quotient}给出 \(E_i/E_{i-1}\) 为 Galois 扩张，且
\[
\operatorname{Gal}(E_i/E_{i-1})\cong H_{i-1}/H_i\cong C_{\ell_i}.
\]
这里所需的正规性是 \(H_i\trianglelefteq H_{i-1}\)，并不要求 \(H_i\trianglelefteq H\)：每一步 \(E_i/E_{i-1}\) 都是循环 Galois 扩张，而 \(E_i/F\) 可以不正规。素数 \(\ell_i\) 整除 \(|H|\)，\(|H|\) 又整除 \(m\)，故 \(\ell_i\mid m\)。因此 \(F\) 及所有 \(E_{i-1}\) 都含 \(\mu_{\ell_i}\)，特征零又使 \(\ell_i\) 在这些域中非零。对每一步应用\cref{b1:thm:kummer-cyclic}，得到
\[
E_i=E_{i-1}(\gamma_i),\qquad \gamma_i^{\ell_i}\in E_{i-1}.
\]
把开头的 \(K\subseteq F\) 接上这些步骤，便得到包含 \(L\) 全部元素、从而包含 \(f\) 全部根的根式塔。
\end{proof}

\begin{corollary}\label{b1:cor:irreducible-one-radical-root}
设 \(K\) 为特征零域，\(f\in K[X]\) 为不可约多项式。若 \(f\) 的一个根属于某条从 \(K\) 出发的根式塔的末端，则 \(f\) 根式可解。
\end{corollary}
\begin{proof}
由\cref{b1:lem:radical-normal-envelope}，可将该根式塔置于有限 Galois 扩张 \(M/K\) 中，且 \(\operatorname{Gal}(M/K)\) 可解。由于 \(f\) 不可约，它与所给根的最小多项式只相差一个非零常数因子。\(M/K\) 正规，故这个最小多项式的全部共轭根都在 \(M\) 中，即 \(f\) 在 \(M\) 中完全分裂。其分裂域 \(L\) 是 \(M\) 的中间域；由\cref{b1:thm:galois-quotient}，\(\operatorname{Gal}(L/K)\) 是 \(\operatorname{Gal}(M/K)\) 的商群，也可解。再由\cref{b1:thm:radical-solvability}，\(f\) 根式可解。
\end{proof}

反向蕴含从根式可解的定义直接得到。因此，对特征零域上的不可约多项式，一个根能否用根式表示与全部根能否用根式表示是同一个存在性问题。不可约条件不能删去：若 \(f\) 可约，已知的一个根只约束它所属的那个不可约因子，不能据此判断其余因子的根式可解性。

\ProseParagraphBreak
可解扩张不一定自身是根式扩张。取本原七次单位根 \(\zeta_7=e^{2\pi i/7}\)，并令
\[
E=\mathbb Q(\zeta_7+\zeta_7^{-1})\subseteq\mathbb R.
\]
由\cref{b1:thm:cyclotomic-galois-group}，\(\operatorname{Gal}(\mathbb Q(\zeta_7)/\mathbb Q)\cong(\mathbb Z/7\mathbb Z)^\times\cong C_6\)，其中 \(3\bmod7\) 的阶为 \(6\)。复共轭的不动域是 \(E\)，它生成的二阶子群正规于 \(C_6\)，所以\cref{b1:thm:galois-quotient}给出 \(E/\mathbb Q\) 为 Galois 扩张，且 \(\operatorname{Gal}(E/\mathbb Q)\cong C_3\)。因此 \(E/\mathbb Q\) 是可解扩张，但以下论证将说明：开方若始终限制在这个实域 \(E\) 内，就不能生成 \(E\)。

\ProseParagraphBreak
若 \(0\neq\beta\in E\)，且对某个整数 \(m\geq2\) 有 \(\beta^m\in\mathbb Q\)，则对任意 \(\sigma\in\operatorname{Gal}(E/\mathbb Q)\)，比值 \(\sigma(\beta)/\beta\) 是实数域 \(E\) 中的 \(m\) 次单位根，因而属于 \(\{1,-1\}\subseteq\mathbb Q\)。这些比值被每个自同构固定，所以与前面根式扩张中的比值计算一样，\(\sigma\mapsto\sigma(\beta)/\beta\) 是到 \(\{1,-1\}\) 的群同态。它的像的阶必须同时整除 \(3\) 和 \(2\)，故只能为平凡群。这说明全部自同构都固定 \(\beta\)，由 Galois 对应得到 \(\beta\in\mathbb Q\)。

\ProseParagraphBreak
若某条根式塔从 \(\mathbb Q\) 出发而以 \(E\) 为末端，则各层都在 \(E\) 内，第一步非平凡扩张须由这样的 \(\beta\in E\setminus\mathbb Q\) 生成，矛盾。所以 \(E/\mathbb Q\) 本身不是根式扩张。另一方面，\(\zeta_7+\zeta_7^{-1}\) 生成 \(E\)，而 \(E/\mathbb Q\) 正规，故它的最小多项式的全部根都在 \(E\) 内，并生成 \(E\)。这个多项式的分裂域便是 \(E\)，由\cref{b1:thm:radical-solvability}，它根式可解。因此 \(E\) 仍包含于某条从 \(\mathbb Q\) 出发的根式塔的末端，只是那个末端域必须严格大于 \(E\)。

\subsection{特征条件与可分性}
\label{b1:subsec:2102:03}

特征零保证每个不可约多项式可分，并允许所取根式指数作为非零域元素使用。\(f\) 本身可以含重因子；删去重数既不改变分裂域，也不改变根式可解性。因此定理并不另要求 \(f\) 的各因子互异。

\ProseParagraphBreak
在特征 \(p\) 中，若只允许指数与 \(p\) 互素的根式，\cref{b1:lem:radical-normal-envelope}的证明仍能构造可解的正规包络。但循环 \(p\) 次扩张由 \(X^p-X-a\) 产生，不能调用\cref{b1:thm:kummer-cyclic}把这一步变为 \(p\) 次开方。另一方面，\(X^p-a\) 的根给出纯不可分扩张；纯不可分扩张的每个元素在基域上的共轭只有自身，所以其基域自同构群必为平凡群。当扩张非平凡时，它不是有限 Galois 扩张，不能用有限 Galois 对应来分析。正特征下需要把 Kummer 步骤、Artin--Schreier 步骤和纯不可分步骤分别处理。

\ProseParagraphBreak
可解群的次正规列仍能指导正特征下的构造：若某个素数阶因子等于特征，使用\cref{b1:thm:artin-schreier-classification}。得到的是允许 Kummer 和 Artin--Schreier 两类步骤的扩张塔，与只允许开方的根式塔应当区分。

% !TeX root = ../../Book1.tex
% 纲要标题：### 21.3 经典方程
% 纲要位置：计划纲要.md，第 752 行。

\section{经典方程}
\label{b1:sec:2103}

二次、三次和四次方程的根式公式并非互不相关的技巧。它们背后的 Galois 群都嵌在可解的低阶对称群中；五次开始，对称群的可解性失效。不过，“一般方程没有通用根式公式”不表示每个具体五次方程都不能用根式求解。

% 纲要内容（仅作写作计划，尚非教材正文）：
%
% * 三次、四次方程的群论解释；从三次根的循环作用理解配对公式，从四根的三种配对构造辅助三次方程并联系 Ferrari 配方
% * 一般五次方程不可根式求解
% * 一个具体不可解五次方程：以 \(x^5-x-1\) 的模素数分解判定 Galois 群，引用第19.3节已证的循环型判据
% * 可解但非交换的 Galois 群：回引三次分裂域的群计算，解释正规子群链、单位根和根式塔
%

\subsection{三次、四次方程的群论解释}
\label{b1:subsec:2103:01}

设基域特征为零。若次数为 \(d\) 的多项式有 \(s\) 个不同的根，则其 Galois 群忠实作用在这 \(s\) 个根组成的集合上，因而嵌入 \(S_s\)，再将 \(S_s\) 看成 \(S_d\) 的子群即可。这里置换的是不同的根；即使按重数列出 \(d\) 个根，也不能据此增加域自同构。式\eqref{b1:eq:small-symmetric-derived}已经给出
\[
S_3\trianglerighteq A_3\trianglerighteq1,
\qquad S_4\trianglerighteq A_4\trianglerighteq V_4\trianglerighteq1
\]
都有交换因子。\cref{b1:thm:solvable-closure}把可解性传给子群，再由\cref{b1:thm:radical-solvability}得到三次、四次方程都根式可解。

\ProseParagraphBreak
可解性判据说明根式表达存在，却没有直接给出求根的步骤。三次方程的 Cardano 公式可以从根的置换中构造出来。将一般三次方程除以首项系数，再平移变量消去二次项，得到 \(y^3+py+q=0\)。设其三个根为 \(r_1,r_2,r_3\)，按重数列出，于是 \(r_1+r_2+r_3=0\)。先加入本原三次单位根 \(\omega\)。我们希望选取根的线性组合，使三轮换的作用成为乘以 \(\omega\) 或 \(\omega^2\)，从而把循环作用化为特征向量上的标量作用。下面两组系数正是为此选择的：
\[
u=\frac{r_1+\omega^2r_2+\omega r_3}{3},\qquad
v=\frac{r_1+\omega r_2+\omega^2r_3}{3}.
\]
由 \(1+\omega+\omega^2=0\)，反过来有
\[
r_1=u+v,\qquad r_2=\omega u+\omega^2v,\qquad
r_3=\omega^2u+\omega v.
\]
在固定 \(\omega\) 的条件下，循环改标 \((r_1,r_2,r_3)\mapsto(r_2,r_3,r_1)\) 使 \((u,v)\mapsto(\omega u,\omega^2v)\)，而交换 \(r_2,r_3\) 使 \(u,v\) 交换。因此 \(u^3,v^3\) 不受三轮换影响，却仍可能互换。这里的改标首先是对根的形式重排，并没有断言某个具体方程一定存在实现它的域自同构。对无重根方程，加入 \(\omega\) 后的 Galois 群中实际出现的置换才按这些公式作用；这个群可以严格小于 \(S_3\)。就上述形式置换而言，商群 \(S_3/A_3\cong C_2\) 记录 \(u^3,v^3\) 的交换，因而先用至多一次开平方将这两个量区分；随后 \(A_3\cong C_3\) 在 \(u,v\) 上分别乘以 \(\omega,\omega^2\)，对应于开立方恢复它们。

\ProseParagraphBreak
Vieta 公式给出 \(r_1+r_2+r_3=0\) 与 \(r_1r_2+r_1r_3+r_2r_3=p\)，所以
\[
r_1^2+r_2^2+r_3^2
=(r_1+r_2+r_3)^2-2(r_1r_2+r_1r_3+r_2r_3)=-2p.
\]
展开 \(u,v\) 的乘积并使用 \(1+\omega+\omega^2=0\)，便得
\[
9uv=r_1^2+r_2^2+r_3^2-(r_1r_2+r_1r_3+r_2r_3)
=-2p-p=-3p,
\]
所以 \(uv=-p/3\)，又有 \(u^3+v^3=r_1^3-3uvr_1=r_1^3+pr_1=-q\)。故 \(u^3,v^3\) 是二次方程
\[
T^2+qT-\frac{p^3}{27}=0
\]
的两个根。这就解释了 Cardano 公式中的平方根为何先于立方根出现。当 \(p=0\) 时直接开立方即可。当 \(p\neq0\) 时，设
\[
\Delta=\frac{q^2}{4}+\frac{p^3}{27},\qquad
U=-\frac q2+\sqrt\Delta,\qquad V=-\frac q2-\sqrt\Delta.
\]
有 \(U+V=-q\)、\(UV=-p^3/27\neq0\)。任取 \(u^3=U\)，再按乘积条件定义 \(v=-p/(3u)\)，则 \(v^3=V\)。于是
\[
(u+v)^3+p(u+v)+q=u^3+v^3+(3uv+p)(u+v)+q=0.
\]
利用已加入的 \(\omega\)，三个根为 \(\omega^j u+\omega^{-j}v\)、\(j=0,1,2\)，重根情形按重数理解。不能独立选取两个立方根后直接相加；\(uv=-p/3\) 是公式成立的必要配对条件。

\ProseParagraphBreak
实系数方程的计算途中也可能出现复数。例如 \(y^3-3y+1=0\) 的 \(\Delta=1/4-1=-3/4\)，选取复平方根后有 \(U=-1/2+i\sqrt3/2\)。因为 \(|U|=1\)，任一立方根 \(u\) 也满足 \(|u|=1\)，配对条件给出 \(v=1/u=\bar u\)。又因 \(\omega^{-j}=\overline{\omega^j}\)，三个表达式 \(\omega^j u+\omega^{-j}v\) 全为实数，计算却经过了复数中间量。这正说明根式可解允许使用复数中间量，并不要求每次开方都留在实域中。

\ProseParagraphBreak
四次方程也能通过一个辅助三次方程归结为二次方程。经平移后，它成为 \(y^4+py^2+qy+r=0\)。若 \(q=0\)，先将 \(y^2\) 作为未知量解二次方程。若 \(q\neq0\)，希望把四次式写成两个平方之差，从而分解为两个二次式。引入待定参数 \(m\)，有
\[
y^4+py^2+qy+r
=(y^2+m)^2-\bigl((2m-p)y^2-qy+m^2-r\bigr).
\]
要使括号中的二次式成为 \(\bigl(sy-q/(2s)\bigr)^2\)，比较二次项与常数项就要求
\[
s^2=2m-p,\qquad m^2-r=\frac{q^2}{4s^2}.
\]
消去 \(s\)，得到关于 \(m\) 的三次方程
\[
4(2m-p)(m^2-r)=q^2
\]
取它的任一根 \(m\)，因为 \(q\neq0\)，有 \(2m-p\neq0\)，可选取非零 \(s\) 满足 \(s^2=2m-p\)。三次方程保证上述常数项条件也成立，于是
\begin{align*}
y^4+py^2+qy+r
&=(y^2+m)^2-\Bigl(sy-\frac q{2s}\Bigr)^2\\
&=\Bigl(y^2-sy+m+\frac q{2s}\Bigr)
  \Bigl(y^2+sy+m-\frac q{2s}\Bigr).
\end{align*}
后面只剩两个二次方程。这说明四次求解为何能够归结为三次求解及连续开平方。

\ProseParagraphBreak
上述配方给出了两个二次因子，还可以从根的配对解释辅助三次方程的来源。设
\[
f(X)=X^4+bX^3+cX^2+dX+e
=\prod_{i=1}^4(X-\alpha_i),
\]
把四个根的指标分成两组，每组两个；不区分组内顺序，也不区分两组的先后。这三种无序划分为 \(12\mid34\)、\(13\mid24\)、\(14\mid23\)，分别对应下面三个配对量：
\[
\begin{aligned}
\beta_1&=\alpha_1\alpha_2+\alpha_3\alpha_4,\\
\beta_2&=\alpha_1\alpha_3+\alpha_2\alpha_4,\\
\beta_3&=\alpha_1\alpha_4+\alpha_2\alpha_3.
\end{aligned}
\]
若 \(e_i\) 表示四个根的第 \(i\) 个初等对称多项式，则 \(e_1=-b\)、\(e_2=c\)、\(e_3=-d\)、\(e_4=e\)。要找一个以三个 \(\beta_j\) 为根的基域多项式，需要把它们的初等对称式改写为 \(e_i\)。其中两两乘积之和展开后，恰含十二个形如 \(\alpha_i^2\alpha_j\alpha_k\)、\(i,j,k\) 互异的单项式，各出现一次；\(e_1e_3\) 包含同样的十二项，并另有四项 \(e_4\)，所以两者相差 \(4e_4\)。三重乘积的八项则可分成两类：
\[
\beta_1\beta_2\beta_3
=e_4\sum_{i=1}^4\alpha_i^2
 +\sum_{1\le i<j<k\le4}(\alpha_i\alpha_j\alpha_k)^2.
\]
第一个和为 \(e_1^2-2e_2\)；第二个和为 \(e_3^2-2e_2e_4\)，因为四个三重乘积的六个两两乘积之和恰为 \(e_2e_4\)。于是得到
\begin{align*}
\beta_1+\beta_2+\beta_3&=e_2=c,\\
\beta_1\beta_2+\beta_1\beta_3+\beta_2\beta_3
&=e_1e_3-4e_4=bd-4e,\\
\beta_1\beta_2\beta_3&=e_1^2e_4+e_3^2-4e_2e_4
=b^2e+d^2-4ce.
\end{align*}
因此它们满足同一个三次辅助方程，其多项式为
\begin{equation}\label{b1:eq:quartic-cubic-resolvent}
\begin{aligned}
R(T)&\coloneqq\prod_{j=1}^3(T-\beta_j)\\
&=T^3-cT^2+(bd-4e)T+4ce-b^2e-d^2.
\end{aligned}
\end{equation}
置换四个根就置换这三个量，这正是\subsecref{b1:subsec:0404:03}中三种二元划分的作用。它的核为
\[
V_4=\{1,(12)(34),(13)(24),(14)(23)\}.
\]
这些置换保持每一种无序划分的整体，允许交换组内的根，也允许交换两组。若四根互异，由
\[
\beta_1-\beta_2=(\alpha_1-\alpha_4)(\alpha_2-\alpha_3)
\]
及另两式可知三个 \(\beta_j\) 也互异。因此保持全部三种划分等价于固定全部三个配对量，\(V_4\) 在这里有了具体的域论含义。仍在四根互异的条件下，原四次方程的 Galois 群 \(G\leq S_4\) 在这些量上的作用核就是 \(G\cap V_4\)。

\ProseParagraphBreak
现在将配对量与配方中的参数对应起来。对前面的 \(y^4+py^2+qy+r\)，有 \(b=0\)、\(c=p\)、\(d=q\)、\(e=r\)，故
\[
R(2m)=4(2m-p)(m^2-r)-q^2.
\]
辅助方程使 \(R(2m)=0\)，而 \(R\) 的根正是三个 \(\beta_j\)，故 \(2m=\beta_j\) 对某个 \(j\) 成立。重新标号，把该配对写为 \(\{\alpha_1,\alpha_2\}\mid\{\alpha_3,\alpha_4\}\)，于是 \(2m=\alpha_1\alpha_2+\alpha_3\alpha_4\)。暂令 \(s_0=\alpha_1+\alpha_2=-(\alpha_3+\alpha_4)\)。四个交叉乘积之和为 \((\alpha_1+\alpha_2)(\alpha_3+\alpha_4)=-s_0^2\)，故 Vieta 公式给出
\[
p=\alpha_1\alpha_2+\alpha_3\alpha_4
  +(\alpha_1+\alpha_2)(\alpha_3+\alpha_4)=2m-s_0^2.
\]
再比较两组根组成的两个二次因子的一次项系数，得到
\[
s_0^2=2m-p,\qquad
q=s_0(\alpha_1\alpha_2-\alpha_3\alpha_4),\qquad
\alpha_1\alpha_2+\alpha_3\alpha_4=2m.
\]
因此 \(s_0=\pm s\)；若 \(s_0=-s\)，交换两组根后便有 \(s_0=s\)。当 \(q\neq0\) 时，\(s\neq0\)，两组根的积于是分别为 \(m+q/(2s)\) 与 \(m-q/(2s)\)，恰好得到上面两个二次因子。辅助三次方程确定配对量，开平方确定两组根的和，最后各解一个二次方程；群作用与求根计算在此描述的是同一过程。

\subsection{一般五次方程不可根式求解}
\label{b1:subsec:2103:02}

下面的“一般”有一个精确含义：把系数当作代数独立的变量，而不是先给它们某组具体数值。

\begin{theorem}[一般方程]\label{b1:thm:generic-equation}
设 \(n\geq1\) 为整数，并设 \(u_1,\ldots,u_n\) 在 \(\mathbb Q\) 上代数独立。多项式
\[
X^n-u_1X^{n-1}+u_2X^{n-2}-\cdots+(-1)^nu_n
\]
在 \(\mathbb Q(u_1,\ldots,u_n)\) 上的 Galois 群是 \(S_n\)。因此当 \(n\geq5\) 时，一般 \(n\) 次方程不根式可解。
\end{theorem}
\begin{proof}
先取代数独立的 \(t_1,\ldots,t_n\)，令 \(E=\mathbb Q(t_1,\ldots,t_n)\)，\(e_j\) 为这些变量的第 \(j\) 个初等对称多项式，并令 \(F=\mathbb Q(e_1,\ldots,e_n)\)。于是
\[
\prod_{i=1}^n(X-t_i)=X^n-e_1X^{n-1}+\cdots+(-1)^ne_n.
\]
每个 \(t_i\) 都是右端多项式的根，因而在 \(F\) 上代数。依次加入 \(t_1,t_2,\ldots,t_n\)，在已知一个根后除去对应线性因子，下一根满足的多项式次数就降低一。所以各步次数分别至多为 \(n,n-1,\ldots,1\)，由\cref{b1:thm:tower-law}得到 \([E:F]\leq n!\)，特别地，扩张有限。各 \(t_i\) 互异，故 \(E/F\) 是右端可分多项式的分裂域，为 Galois 扩张。

变量的任意置换保持多项式环的加法与乘法，并延拓为分式域 \(E\) 的自同构。每个 \(e_j\) 都是对称多项式，所以这些置换固定全部 \(e_j\)，也就逐点固定它们生成的系数域 \(F\)。不同置换在某个 \(t_i\) 上取值不同，给出 \(n!\) 个不同的 \(F\) 自同构。\cref{b1:thm:separable-embeddings}的嵌入数上界于是给出
\[
n!\leq|\operatorname{Aut}_F(E)|\leq[E:F]\leq n!.
\]
全部等号成立，所有自同构恰为这些变量置换，因此 Galois 群就是 \(S_n\)。

还需证明这些初等对称多项式代数独立。若 \(H\in\mathbb Q[U_1,\ldots,U_n]\) 满足 \(H(e_1,\ldots,e_n)=0\)，因 \(t_1,\ldots,t_n\) 代数独立，这就是关于 \(t_i\) 的多项式恒等式。\cref{b1:thm:fundamental-symmetric-polynomials}的唯一性具体说：一个对称多项式写成初等对称多项式的多项式表达时，所用的表达式唯一。按 \(t_1>\cdots>t_n\) 的词典序，\(e_1^{a_1}\cdots e_n^{a_n}\) 的首项指数依次为 \(a_1+\cdots+a_n,a_2+\cdots+a_n,\ldots,a_n\)，依次作差便恢复各 \(a_i\)，所以不同单项式的首项也不同。现在零多项式同时由 \(H\) 与零表达，所以 \(H=0\)，证明了 \(e_1,\ldots,e_n\) 的代数独立性。由 \(u_i\mapsto e_i\) 得到系数域 \(\mathbb Q(u_1,\ldots,u_n)\cong F\) 的同构，并把题述多项式送到上面已经考察的多项式，因此两者的 Galois 群同构。

当 \(n\geq5\) 时，\cref{b1:cor:symmetric-nonsolvable}说明 \(S_n\) 不可解；\cref{b1:thm:radical-solvability}遂排除根式求解。
\end{proof}

这就是\term[Abel-Ruffini-dingli]{Abel--Ruffini 定理}[Abel--Ruffini theorem]的通用方程形式。这里的“一般”始终指系数域 \(\mathbb Q(u_1,\ldots,u_n)\) 上那一个以代数独立变量为系数的多项式，不是关于随机选取数值系数的概率陈述。它排除的是以这些独立系数为起点的通用根式表达，并不排除特殊方程。例如 \(X^5-1\) 的分裂域是分圆域，Galois 群 \((\mathbb Z/5\mathbb Z)^\times\cong C_4\) 可解，因而它的全部根都能用根式表示。


\subsection{一个不可根式求解的五次方程}
\label{b1:subsec:2103:04}

一般方程的结论不能直接判定某一个给定方程。现在考察序言中的 \(X^5-X-1\)，用\cref{b1:lem:reduction-cycle-type}把模素数分解得到的循环型与群的传递性结合起来，确定它的 Galois 群。


\begin{proposition}\label{b1:prop:quintic-x5-x-1}
多项式 \(f=X^5-X-1\) 在 \(\mathbb Q\) 上的 Galois 群为 \(S_5\)，因而不根式可解。
\end{proposition}
\begin{proof}
先证明不可约性。五次多项式若可约，就至少有两个不可约因子；若每个因子的次数都不小于三，总次数至少为六，矛盾。因此只须排除一次与不可约二次因子。模三时，\(f\) 在 \(0,1,2\) 处的值都为二，没有一次因子。\(\mathbb F_3\) 上的首一不可约二次多项式只有
\(X^2+1\)、\(X^2+X+2\)、\(X^2+2X+2\)：对常数项非零的六个首一二次式逐一检查是否在 \(0,1,2\) 处有根，即得此表。用它们分别除 \(\bar f\)，余式依次为
\[
2,\qquad X+2,\qquad X+2.
\]
所以也没有不可约二次因子，\(\bar f\) 不可约；\cref{b1:prop:reduction-irreducible}说明 \(f\) 在 \(\mathbb Q\) 上不可约。

模二时，直接相乘可核验
\[
\bar f=(X^2+X+1)(X^3+X^2+1).
\]
两个因子在 \(\mathbb F_2\) 中均无根，次数分别为二和三，因而均不可约且互异，所以约化无重因子。有限域上的不可约多项式都可分，故约化也无重根，满足\cref{b1:lem:reduction-cycle-type}的条件。该引理给出群中的一个 \((2)(3)\) 型置换 \(\tau=(a\,b)(d\,e\,h)\)。这两个循环不相交，彼此交换，故
\[
\tau^3=(a\,b)^3(d\,e\,h)^3=(a\,b),
\]
得到群中的一个换位。

在分裂域中，若 \(a,b\) 是同一不可约因子的两个根，\cref{b1:lem:simple-embedding-extension}对基域恒等映射给出 \(K(a)\to K(b)\)、\(a\mapsto b\) 的同构；\cref{b1:thm:embedding-extension}延拓这个同构，正规性使得到的嵌入成为分裂域的自同构。因此一个不可约因子的全部根恰组成一个轨道。现在 \(f\) 不可约，故五个根只有一个轨道，群的作用传递。再由\cref{b1:thm:orbit-stabilizer}，五整除群阶。\cref{b1:thm:cauchy}给出五阶元素，在五个点上它必为五轮换 \(c\)。

将五个点按 \(c\) 标为 \(\mathbb Z/5\mathbb Z\)，使 \(c\) 的作用为加一。设刚找到的换位交换 \(i,j\)，记非零差为 \(d=j-i\)。用 \(c\) 的各次幂共轭这个换位，就得到所有 \((k\,\,k+d)\)。由于五为素数，从任一点反复加 \(d\)，依次经过全部五个点；所以任意两点都可用这些相邻关系连成一条不重复顶点的路径。设路径为 \(v_0,\ldots,v_r\)，它的各相邻换位都在群中。由第一个换位开始，若已经得到 \((v_0\,v_j)\)，就有
\[
(v_0\,v_j)(v_j\,v_{j+1})(v_0\,v_j)=(v_0\,v_{j+1})
\qquad(1\le j<r).
\]
沿路径归纳，得到两端的换位 \((v_0\,v_r)\)。任意两点均可作两端，故群包含全部换位；换位生成 \(S_5\)，于是原 Galois 群等于 \(S_5\)。由\cref{b1:cor:symmetric-nonsolvable}及\cref{b1:thm:radical-solvability}，方程不根式可解。
\end{proof}
这个具体判定解释了序言中 \(X^5-1\) 与 \(X^5-X-1\) 的差别：前者有交换的分圆 Galois 群，后者的群为不可解的 \(S_5\)。次数同为五，并不意味着两者有相同的根式可解性。

\ProseParagraphBreak
对可约多项式，能求出一个根也不意味着全部根可用根式求出。例如 \((X-1)(X^5-X-1)\) 有有理根 \(1\)，但它的五次因子已由上面的命题证明不根式可解。

\subsection{可解但非交换的 Galois 群}
\label{b1:subsec:2103:03}

回到 \(X^3-2\)。令 \(a=\sqrt[3]{2}\)、\(\zeta=\zeta_3\)、\(L=\mathbb Q(a,\zeta)\)。\cref{b1:prop:cubic-splitting-field-correspondence}已确定 \([L:\mathbb Q]=6\)、\(\operatorname{Gal}(L/\mathbb Q)\cong S_3\) 及其全部中间域。这里利用这些结论，解释根式求解怎样反映在子群与域塔上。

\ProseParagraphBreak
由 Galois 对应，正规子群链和相应的域塔为
\[
S_3\trianglerighteq A_3\trianglerighteq1,
\qquad
\mathbb Q\subseteq\mathbb Q(\zeta)\subseteq L.
\]
两步扩张的 Galois 群分别是 \(S_3/A_3\cong C_2\) 与 \(A_3\cong C_3\)。第一步通过
\[
\zeta=\frac{-1+\sqrt{-3}}2,
\qquad \mathbb Q(\zeta)=\mathbb Q(\sqrt{-3})
\]
开平方完成，反向关系 \(\sqrt{-3}=2\zeta+1\) 也说明两域相等。这一步使用复平方根，已离开实域，仍属于允许的根式运算。第二步加入 \(a\)，其立方 \(a^3=2\) 已在基域中。于是这条对应于正规子群链的域塔本身就是一条根式塔。

\ProseParagraphBreak
先加入单位根，使第二步也能由 Galois 理论描述：在 \(\mathbb Q(\zeta)\) 上，加入 \(a\) 便同时包含 \(a,\zeta a,\zeta^2a\) 三个共轭根，循环自同构把 \(a\) 送到 \(\zeta a\)。若先走 \(\mathbb Q\subseteq\mathbb Q(a)\)，虽然已经通过开立方取得一个根，这一步却不正规，因为另外两个根不在实域 \(\mathbb Q(a)\) 中。

\ProseParagraphBreak
若 \(\rho(a)=\zeta a\)、\(\rho(\zeta)=\zeta\)，而 \(\kappa\) 为复共轭，则 \(\kappa\rho(a)=\zeta^2a\)、\(\rho\kappa(a)=\zeta a\)，两者不同，故全群 \(S_3\) 非交换。但上述两层商群 \(S_3/A_3\cong C_2\) 与 \(A_3\cong C_3\) 都是交换群，所以 \(S_3\) 仍然可解。交换群本身总可解，而这个例子说明可解群可以非交换；根式求解所需要的是逐层出现交换的商群，不要求整个 Galois 群交换。

% !TeX root = ../../Book1.tex
% 纲要标题：### 21.4 尺规作图
% 纲要位置：计划纲要.md，第 724 行。

\section{尺规作图}
\label{b1:sec:2104}

尺规作图限制了允许执行的操作：直尺只能连接已知点，圆规只能以已知点为圆心、以已知线段为半径画圆。把点写成坐标以后，直线与圆的交点只要求解一次或二次方程。因此，这种几何限制可以转化为扩张次数的限制。

% 纲要内容（仅作写作计划，尚非教材正文）：
%
% * 二次扩张塔与规矩数
% * 倍立方体、一般三等分角与正多边形
% * 化圆为方与 π 的超越性；超越性论证给出证明概要
%
% ---
%

\subsection{二次扩张塔与规矩数}
\label{b1:subsec:2104:01}

\begin{definition}\label{b1:def:constructible-number}
从平面上的 \((0,0)\)、\((1,0)\) 出发，允许有限次画过两个不同已知点的直线、以已知点为圆心且以两个已知点间距离为半径的圆，并把任意两条所画曲线的孤立交点加入已知点。重合曲线不由此指定新的点。如果一个点能经过有限次这种操作得到，则称它为\term[keguitu-dian]{尺规可作图点}[constructible point]。给定实数 \(a\)，如果 \((a,0)\) 是尺规可作图点，则称 \(a\) 为\term[guiju-shu]{规矩数}[constructible number]，也称可构造数、可作图数。\footnote{北京交通大学《近世代数》课程使用“规矩数”\cite[第 23 章第 1 节“几何作图简史、规矩数”]{BJTU2026ModernAlgebraCourse}；“可构造数”见席南华《基础代数》第一卷\cite[p. 125]{Xi2016BasicAlgebraI}，该书从复平面讨论可构造点。}\termidx[kegouzao-shu]{可构造数}\termidx[kezuotu-shu]{可作图数}
\end{definition}

标准的作垂线、平行线与转移线段长度均由上述操作实现。因此一个点可作图当且仅当它的两个坐标都是规矩数：从点向坐标轴作垂线得到坐标，反过来在两轴上标出坐标再作平行线，交点即为所需点。

\begin{theorem}\label{b1:thm:constructible-quadratic-tower}
实数 \(a\) 可作图，当且仅当存在全部位于 \(\mathbb R\) 中的域塔
\[
\mathbb Q=K_0\subseteq K_1\subseteq\cdots\subseteq K_r,
\qquad [K_i:K_{i-1}]=2,
\]
使 \(a\in K_r\)。允许 \(r=0\)。特别地，规矩数在 \(\mathbb Q\) 上的次数必为 \(2\) 的幂。
\end{theorem}
\begin{proof}
若当前全部已知坐标在实域 \(F\) 中，过已知点的直线方程系数属于 \(F\)。以已知点为圆心、已知线段为半径的圆方程是
\[
(x-u)^2+(y-v)^2=r^2,
\qquad u,v,r^2\in F.
\]
两直线的孤立交点由一次方程求出，坐标仍在 \(F\)。一条直线与一个圆联立，消去一个变量便得到至多二次方程；求出该变量后，另一个变量是它的 \(F\) 线性表达式。两圆相减得到一次方程，故也归结为前一种情况。重合的曲线不指定新的孤立交点；无交点时没有可执行步骤。因此每次新添交点至多需要一个实二次扩张。沿有限作图过程依次加入坐标，删去次数为 \(1\) 的步骤，就得到所需塔。

反过来，可作图的有向线段长度对加、减封闭，乘、除由相似三角形实现。例如连结 \((1,0)\) 与 \((0,b)\)，过 \((a,0)\) 作平行线，其纵轴截距为 \(ab\)。若 \(a\neq0\)，改为连结 \((a,0)\) 与 \((0,b)\)，过 \((1,0)\) 作平行线，纵轴截距便是 \(b/a\)。这些操作连同符号方向使规矩数构成一个实子域。

若 \(u>0\) 已可作图，在横轴上取 \(A=(0,0)\)、\(C=(1,0)\)、\(B=(1+u,0)\)，画以 \(AB\) 为直径的半圆，并过 \(C\) 作垂线，与上半圆交于 \(P=(1,h)\)，如\cref{b1:fig:square-root-semicircle}所示。圆心为 \(AB\) 的中点，半径为 \(AB\) 的一半，故
\[
\left(1-\frac{1+u}{2}\right)^2+h^2
=\left(\frac{1+u}{2}\right)^2,
\qquad h^2=u.
\]
取上半圆使 \(h>0\)，便作出了 \(\sqrt u\)。

% !TeX root = ../Book1.tex
\begin{figure}[htbp]
\centering
\begin{tikzpicture}[scale=1.35]
  % The proportions illustrate u=9/4; the construction works for every u>0.
  \pgfmathsetmacro{\diagramu}{2.25}
  \pgfmathsetmacro{\diagramr}{(1+\diagramu)/2}
  \pgfmathsetmacro{\diagramh}{sqrt(\diagramu)}
  \coordinate (A) at (0,0);
  \coordinate (C) at (1,0);
  \coordinate (B) at ({1+\diagramu},0);
  \coordinate (O) at (\diagramr,0);
  \coordinate (P) at (1,\diagramh);
  \draw[thick] (A)--(B);
  \draw[thick,color=AlgebraBlue] (A) arc[start angle=180,end angle=0,radius=\diagramr];
  \draw[thick,color=AlgebraBlue] (C)--(P);
  \draw[gray,dashed] (O)--(P);
  \draw (1,0.14)--(1.14,0.14)--(1.14,0);
  \foreach \point in {A,B,C,O,P} {\fill (\point) circle (1.2pt);}
  \node[below left] at (A) {\(A\)};
  \node[below] at (C) {\(C\)};
  \node[below right] at (B) {\(B\)};
  \node[below right] at (O) {\(O\)};
  \node[above left] at (P) {\(P\)};
  \node[below] at (0.5,-0.29) {\(1\)};
  \node[below] at ({1+\diagramu/2},-0.29) {\(u\)};
  \node[left,color=AlgebraBlue] at (0.92,{\diagramh/2}) {\(h\)};
\end{tikzpicture}
\caption[半圆开平方作图]{半圆开平方作图。\(AC=1\)、\(CB=u>0\)，\(O\) 为 \(AB\) 的中点，过 \(C\) 的垂线与上半圆交于 \(P\)。高度 \(CP=h\) 满足 \(h^2=u\)。图形比例仅作示意。}
\label{b1:fig:square-root-semicircle}
\end{figure}


最后，对任意实二次扩张 \(F(\beta)/F\)，把 \(\beta\) 的首一最小多项式写成 \(X^2+bX+c\)，配方得到 \((2\beta+b)^2=b^2-4c\)。判别式属于 \(F\) 且为正，所以 \(\beta=(-b\pm\sqrt{b^2-4c})/2\) 可由域运算和开正平方根得到。对塔归纳，末端全部元素都可作图。由\cref{b1:thm:tower-law}，\([K_r:\mathbb Q]=2^r\)，而 \(\bigl[\mathbb Q(a):\mathbb Q\bigr]\) 整除它，故为 \(2\) 的幂。
\end{proof}

若从更多给定点出发，应把 \(\mathbb Q\) 换为这些点坐标生成的域，证明逐字适用。次数为 \(2\) 的幂只是一个必要条件；完整判据要求存在二次扩张塔，不能省掉“塔”而只检查单个次数。

\subsection{倍立方体、三等分角与正多边形}
\label{b1:subsec:2104:02}

倍立方体要求从单位棱长作出棱长 \(\sqrt[3]{2}\)。\cref{b1:thm:eisenstein}对素元 \(2\) 给出 \(X^3-2\) 在 \(\mathbb Q\) 上不可约，故该数次数为 \(3\)，由\cref{b1:thm:constructible-quadratic-tower}不可作图。

\ProseParagraphBreak
一般三等分角也不可能，因为已经可作出的 \(60^\circ\) 角不能被尺规三等分。若 \(20^\circ\) 可作图，则 \(x=2\cos20^\circ\) 可作图；三倍角恒等式给出 \(x^3-3x-1=0\)。由\cref{b1:prop:rational-root}，这个首一整系数三次多项式的有理根只能是 \(\pm1\)，代入都非零。三次多项式可约必有一次因子，因此它不可约，\(x\) 的次数为 \(3\)，矛盾。这里排除的是一种对任意角都适用的作图方法；例如 \(90^\circ\) 可以三等分，并不与结论冲突。

\ProseParagraphBreak
正多边形问题的完整回答是\term[Gauss-Wantzel-dingli]{Gauss--Wantzel 定理}[Gauss--Wantzel theorem]。\footnote{旺策尔（Pierre Laurent Wantzel，1814-1848），法国数学家，证明若干古典尺规作图问题不可能完成。}

\begin{theorem}[Gauss--Wantzel]\label{b1:thm:constructible-polygon}
设 \(n\geq3\) 为整数。正 \(n\) 边形可尺规作图，当且仅当 \(\varphi(n)\) 为 \(2\) 的幂；等价地，
\[
n=2^a p_1\cdots p_s,
\]
其中 \(a,s\geq0\) 为整数，各 \(p_i\) 是两两不同的奇素数，并且每个 \(p_i\) 都可写成 \(2^{2^{r_i}}+1\)，其中 \(r_i\geq0\) 为整数；这样的素数称为\term[Fermat-sushu]{Fermat 素数}[Fermat prime]。\footnote{费马（Pierre de Fermat，1601？-1665），法国数学家，研究数论与解析几何；其出生年份有不同记载。}
\end{theorem}
\begin{proof}
正 \(n\) 边形可作图等价于 \(\zeta_n\) 的实部和虚部可作图。若可作图，把这两个坐标所在的实二次塔合并，再加入 \(i\)，得到一个次数为 \(2\) 的幂的域，其中含有 \(\zeta_n\)。由\cref{b1:thm:cyclotomic-irreducible}及塔公式，\(\varphi(n)=\bigl[\mathbb Q(\zeta_n):\mathbb Q\bigr]\) 也是 \(2\) 的幂。

反过来，若 \(\varphi(n)\) 为 \(2\) 的幂，\cref{b1:thm:cyclotomic-galois-group}说明分圆域的 Galois 群是有限交换 \(2\) 群。复共轭子群的固定域是实子域 \(F=\mathbb Q(\zeta_n+\zeta_n^{-1})\)：这里 \(n\geq3\)，故 \(\zeta_n\) 不实，且它在 \(F\) 上满足 \(X^2-(\zeta_n+\zeta_n^{-1})X+1=0\)。复共轭子群在交换群中正规，所以\cref{b1:thm:galois-quotient}说明 \(F/\mathbb Q\) 也是 Galois 扩张，其群是有限交换 \(2\) 群。这样的群有一条各步指数均为 \(2\) 的子群列：对非平凡有限交换 \(2\) 群取极大真子群，商为交换单群，由\cref{b1:prop:abelian-simple}为 \(C_2\)，再对阶较小的子群归纳。由\cref{b1:thm:finite-galois-correspondence}，\(F\) 因而位于一条实二次扩张塔的末端。于是 \(\cos(2\pi/n)\) 可作图，\(\sin(2\pi/n)\) 也可由 \(\sqrt{1-\cos^2(2\pi/n)}\) 作出，正 \(n\) 边形便可作图。

最后设 \(n=2^a\prod p^{e_p}\)，其中 \(p\) 为奇素数。由模素数幂计数及中国剩余定理，
\[
\varphi(n)=\varphi(2^a)\prod_{p\mid n,\ p\text{ 奇}}p^{e_p-1}(p-1).
\]
它为 \(2\) 的幂当且仅当每个奇素数的指数 \(e_p=1\)，且 \(p-1\) 是 \(2\) 的幂。若 \(p=2^m+1\) 是素数，写 \(m=2^r u\)，其中 \(u\) 为奇数。若 \(u>1\)，则 \(z+1\) 整除 \(z^u+1\)，取 \(z=2^{2^r}\) 会给出 \(p\) 的非平凡因子。因此 \(u=1\)，即 \(m=2^r\)。反向代入上式立即得到 \(\varphi(n)\) 为 \(2\) 的幂。
\end{proof}

例如正 \(3,5,17\) 边形及正 \(15\) 边形可作图，正 \(7\) 边形与正 \(9\) 边形不可作图，因为相应次数为 \(6\)。每个给定的 \(n\) 都可以按其素因数分解单独检查。

\subsection{化圆为方与圆周率的超越性}
\label{b1:subsec:2104:03}

化圆为方要求作出与单位圆等面积的正方形，其边长必须为 \(\sqrt\pi\)。为排除这种长度，需要研究指数函数在代数数处的取值。下述结果是\term[Lindemann-Weierstrass-dingli]{Lindemann--Weierstrass 定理}[Lindemann--Weierstrass theorem]的单变量推论\footnote{林德曼（Carl Louis Ferdinand von Lindemann，1852-1939），德国数学家，证明圆周率的超越性，从而解决化圆为方的不可能性。}；其证明方法将整数的离散性与指数积分的估计结合起来\footnote{魏尔斯特拉斯（Karl Theodor Wilhelm Weierstrass，1815-1897），德国数学家，发展复函数论与分析的严格基础。}。

\begin{theorem}[Lindemann--Weierstrass 的单变量推论]\label{b1:thm:lindemann-weierstrass-special-case}
若 \(\alpha\) 是非零代数数，则 \(e^\alpha\) 是超越数。
\end{theorem}
\begin{proofsketch}
假设 \(e^\alpha=\beta\) 为代数数。取非零多项式 \(P(T)\in\mathbb Z[T]\)，使 \(P(\beta)=0\) 且 \(P(0)\ne0\)，并记 \(\alpha_1=\alpha,\ldots,\alpha_r\) 为 \(\alpha\) 的全部共轭。令
\[
\Lambda=\mathbb Z\alpha_1+\cdots+\mathbb Z\alpha_r\subseteq(\mathbb C,+).
\]
这是有限生成无挠交换群，\cref{b1:thm:fga-structure}保证它有一组整数基 \(\lambda_1,\ldots,\lambda_s\)：每个 \(\gamma\in\Lambda\) 唯一写成 \(\sum_j a_j\lambda_j\)，其中 \(a_j\in\mathbb Z\)。其整群代数 \(\mathbb Z[\Lambda]\) 由形式有限和 \(\sum_\gamma b_\gamma E^\gamma\) 组成，乘法规定为 \(E^\gamma E^\eta=E^{\gamma+\eta}\)。整数基给出同构
\[
\mathbb Z[\Lambda]\cong\mathbb Z[T_1^{\pm1},\ldots,T_s^{\pm1}],
\qquad E^{\sum_j a_j\lambda_j}\longmapsto\prod_j T_j^{a_j}.
\]
右端是 Laurent 多项式环，因而是整环。在这个形式环中构造
\[
Q=\prod_{i=1}^r P(E^{\alpha_i}).
\]
每个 \(\alpha_i\ne0\)，故 \(P(E^{\alpha_i})\) 中不同次数对应不同形式指数，是非零元素；整环性质于是给出 \(Q\ne0\)。可写 \(Q=\sum_\gamma b_\gamma E^\gamma\)，其中 \(b_\gamma\in\mathbb Z\)。代数共轭置换保留 \(Q\)。现在使用求值同态
\[
\operatorname{ev}:\mathbb Z[\Lambda]\longrightarrow\mathbb C,
\qquad E^\gamma\longmapsto e^\gamma,
\]
在 Laurent 多项式坐标中即 \(T_j\mapsto e^{\lambda_j}\)。由 \(P(e^\alpha)=0\)，\(Q\) 的求值为零。这一步并未假定求值同态单射：形式环中的非零元素是否可能求值为零，正是这里要排除的情形。

在形式环中令 \(Q^-=\sum_\gamma b_\gamma E^{-\gamma}\)，即把 \(Q\) 的各指数取负，系数不变。上标 \(-\) 表示这个操作，既不表示乘法逆，也不表示复共轭。乘积 \(QQ^-\) 的常数项为 \(\sum_\gamma b_\gamma^2>0\)，因为一项的指数与另一项的负指数相消，恰好要求原来的两个指数相等。\(QQ^-\) 仍在代数共轭下保持，求值仍为零。合并相同指数，得到
\[
c_0+\sum_{j=1}^n c_j e^{z_j}=0,
\]
其中 \(c_j\) 为整数，\(c_0>0\)，各 \(z_j\) 为互异的非零代数数；指数集合在共轭下保持，且共轭置换保留相应系数。这里共轭只作用于形式表达式的指数；取通常的指数值是随后单独进行的代入。

余下的目标是构造一族不全为零的代数数 \(S_{m,i}\)，使它们乘以一个可控制的整数因子后成为代数整数，同时使全部代数共轭都足够小。非零代数整数的范数是非零整数，绝对值至少为 \(1\)；若所有共轭的乘积反而小于 \(1\)，便得到矛盾。这里约束的是范数，不是某一个共轭的绝对值。

令 \(z_0=0\)、\(f(z)=\prod_{j=0}^n(z-z_j)\)。在各 \(z_j\) 处安排高重零点，可以使许多低阶导数消失；除以 \((m-1)!\) 后，剩余导数中的阶乘比仍为整数，便于控制整性。为把这些导数与一个小积分联系起来，取全部导数之和，使相邻导数相减时消项。具体地，对正整数 \(m\) 及 \(0\le i\le n\)，记
\[
P_{m,i}(z)=\frac{f(z)^m}{z-z_i},\qquad
d_m=(n+1)m-1,\qquad
H_{m,i}(z)=\frac1{(m-1)!}\sum_{r=0}^{d_m}P_{m,i}^{(r)}(z).
\]
这里 \(P_{m,i}\) 是次数为 \(d_m\) 的多项式；它在 \(z_i\) 处有 \(m-1\) 重零，在其余 \(z_k\) 处有 \(m\) 重零。除去一个因子使这些多项式的次数小于 \((n+1)m\)，同时保留后面证明非奇异性所需的重零点。定义端点矩阵 \(A_m=(H_{m,i}(z_k))_{0\le i,k\le n}\)，并沿从 \(0\) 到 \(z_k\) 的线段取积分
\[
\varepsilon_{m,i}(z_k)
=\frac1{(m-1)!}\int_0^{z_k}e^{-z}P_{m,i}(z)\,\mathrm{d}z.
\]
这里的线段积分可用实参数直接表示：将 \(z=tz_k\)、\(0\le t\le1\) 代入，有
\[
\int_0^{z_k}F(z)\,\mathrm{d}z
=z_k\int_0^1F(tz_k)\,\mathrm{d}t,
\]
右端按实部、虚部分别作普通实积分；\(z_k=0\) 时两端均为零。因此积分的绝对值至多为 \(\lvert z_k\rvert\max_{0\le t\le1}\lvert F(tz_k)\rvert\)。
由于 \(H_{m,i}-H_{m,i}'=P_{m,i}/(m-1)!\)，分部积分给出可逐项核对的端点等式
\[
\varepsilon_{m,i}(z_k)
=H_{m,i}(0)-e^{-z_k}H_{m,i}(z_k).
\]

矩阵 \(A_m\) 非奇异，这一点不需要预设其行列式的精确公式。设 \(\sum_i\lambda_iH_{m,i}(z_k)=0\) 对全部 \(k\) 成立，令 \(H=\sum_i\lambda_iH_{m,i}\) 及 \(G=H-H'\)。由定义，\(G\) 被 \(f^{m-1}\) 整除，故在每个 \(z_k\) 处至少有 \(m-1\) 重零。利用 \(H'=H-G\) 逐次求导，并从 \(H(z_k)=0\) 出发归纳，可知 \(H\) 在每个 \(z_k\) 处至少有 \(m\) 重零。可是 \(\deg H\le d_m<(n+1)m\)，所以 \(H=0\)。再由
\[
G(z)=\frac{f(z)^{m-1}}{(m-1)!}
\sum_{i=0}^n\lambda_i\prod_{j\ne i}(z-z_j)
\]
可将括号内的多项式依次在 \(z_i\) 处取值：由于各 \(z_j\) 不同，每次只有第 \(i\) 项不为零，故全部 \(\lambda_i=0\)。

令 \(c=(c_0,\ldots,c_n)^{\mathsf T}\)，并设 \(S_{m,i}=(A_mc)_i=\sum_k c_kH_{m,i}(z_k)\)。由于 \(c_0>0\) 且 \(A_m\) 非奇异，这些代数数不全为零。把端点等式乘以 \(c_ke^{z_k}\) 后求和，由 \(\sum_k c_ke^{z_k}=0\) 得
\[
S_{m,i}=-\sum_{k=0}^n c_ke^{z_k}\varepsilon_{m,i}(z_k).
\]
线段只有有限条；其上 \(e^{-z}\)、\(f(z)\) 与各多项式 \(f(z)/(z-z_i)\) 有共同的有限上界。写 \(P_{m,i}=f^{m-1}\bigl(f/(z-z_i)\bigr)\)，并使用上面的实参数积分界，便得到与 \(m,i\) 无关的 \(C,B>0\)，使 \(\lvert S_{m,i}\rvert\le CB^m/(m-1)!\)。

取正整数 \(q\)，使每个 \(qz_j\) 都是代数整数。\(P_{m,i}\) 在 \(z_k\) 处至少有 \(m-1\) 重零；展开 \(P_{m,i}(z_k+t)\) 后，\(t^r\) 的系数是含 \(d_m-r\) 个差 \(z_k-z_j\) 的整数系数多项式，乘以 \(q^{d_m-r}\) 即成为代数整数。对非零导数项，\(P_{m,i}^{(r)}(z_k)/(m-1)!\) 中的阶乘比 \(r!/(m-1)!\) 是整数。因此取固定的 \(D=q^{n+1}\)，则 \(D^mH_{m,i}(z_k)\)，从而 \(D^mS_{m,i}\)，均为代数整数。

取包含全部 \(z_j\) 的固定 Galois 数域 \(M\)。指数 \(z_j\) 及对应系数 \(c_j\) 在代数共轭下同步置换；若 \(\sigma(z_i)=z_j\)，则 \(\sigma\) 把 \(P_{m,i}\)、\(H_{m,i}\) 分别变为 \(P_{m,j}\)、\(H_{m,j}\)，从而 \(\sigma(S_{m,i})=S_{m,j}\)。故上面的积分界同时控制全部共轭。若 \(D^mS_{m,i}\ne0\)，其范数 \(N_{M/\mathbb Q}(D^mS_{m,i})\) 是非零整数；但其绝对值至多为 \((CD^mB^m/(m-1)!)^{[M:\mathbb Q]}\)，右边随 \(m\) 趋于零。故充分大的 \(m\) 下全部 \(S_{m,i}=0\)。这与它们不全为零矛盾。

原文\cite[\S 1, pp. 214--216]{Lindemann1882Pi}使用 Hermite 积分\footnote{埃尔米特（Charles Hermite，1822-1901），法国数学家，证明自然常数 \(e\) 的超越性，并研究椭圆函数。}及另行构造的系数矩阵；这里直接证明端点矩阵的非奇异性，不沿用原文系数矩阵的精确行列式等式。共轭对称与积分估计在原文\cite[\S\S 2--4, pp. 216--225]{Lindemann1882Pi}中用于推出单变量结论。这个证明概要使用复指数函数、线段上的积分估计和代数整数范数，不能仅由 Galois 对应替代。
\end{proofsketch}

一般定理断言，互异代数数的指数在代数数域上线性无关。取指数 \(0,\alpha\)，若 \(\beta=e^\alpha\) 为代数数，就有代数系数关系 \(-\beta e^0+e^\alpha=0\)，因此上述结论确是它的单变量推论。

\begin{corollary}\label{b1:cor:pi-transcendental}
圆周率 \(\pi\) 是超越数，因而不能用尺规将圆化为等面积的正方形。
\end{corollary}
\begin{proof}
若 \(\pi\) 代数，则非零数 \(i\pi\) 也代数，由\cref{b1:thm:lindemann-weierstrass-special-case}，\(e^{i\pi}\) 应当超越。这与 Euler 恒等式 \(e^{i\pi}=-1\) 矛盾，故 \(\pi\) 超越。

若 \(\sqrt\pi\) 是规矩数，由\cref{b1:thm:constructible-quadratic-tower}，它在 \(\mathbb Q\) 上代数；平方后 \(\pi\) 也应代数，仍然矛盾。因此化圆为方不可能。
\end{proof}
在化圆为方问题中，作图判据先给出可作图数必须代数这一必要条件；超越性论证另使用复指数函数及积分估计。Euler 恒等式也属于复指数函数的性质，因此化圆为方的不可作图性同时包含代数与分析的论证。


\begin{chaptersummary}
在特征零基域中添入所需单位根后，素数阶循环扩张可以由一个根式生成。在 \(\mathbb Q\) 上，分圆多项式的不可约性给出分圆域的次数与完整 Galois 群；在一般基域上，单位根扩张的群是相应单位群的子群。特征零下，求根问题的障碍由可解群刻画：一方面用正规包络把根式塔置于可解 Galois 扩张中，另一方面沿素数阶循环商逐层加入根式。\(S_3\)、\(S_4\) 可解而 \(S_n\)、\(n\geq5\)，不可解，解释了通用公式在五次处发生的变化。

\ProseParagraphBreak
尺规作图只允许实二次步骤，所以比一般根式表达更受限制。作图判据的核心是一条二次扩张塔；可作图实代数数的最小多项式还必须有 \(2\)-群作为分裂域的 Galois 群，单看该数的次数为 \(2\) 的幂并不够。分圆域的交换性又使正多边形的判据简化为 \(\varphi(n)\) 是否为 \(2\) 的幂。对化圆为方，域论只能说明可作图长度必须代数；排除 \(\sqrt\pi\) 还要使用独立的超越性结果。
\end{chaptersummary}

\BookExercises

\begin{exercise}\label{b1:ex:21-1}
计算 \(\Phi_9(X)\)、\(\Phi_{12}(X)\)，并核对它们的次数。
\end{exercise}
\begin{solution}
由\eqref{b1:eq:cyclotomic-factorization}，
\[
\Phi_9(X)=\frac{X^9-1}{X^3-1}=X^6+X^3+1.
\]
又因 \(X^{12}-1=(X^6-1)(X^6+1)\)，在除去阶整除 \(6\) 的根后，\(X^6+1\) 收集阶为 \(4\) 或 \(12\) 的根，所以
\[
\Phi_{12}(X)=\frac{X^6+1}{X^2+1}=X^4-X^2+1.
\]
相应次数为 \(6=\varphi(9)\)、\(4=\varphi(12)\)，两多项式由\cref{b1:thm:cyclotomic-irreducible}均不可约。
\end{solution}

\begin{exercise}\label{b1:ex:21-2}
设 \(n>1\) 为奇数。证明 \(\Phi_{2n}(X)=\Phi_n(-X)\) 及 \(\mathbb Q(\zeta_{2n})=\mathbb Q(\zeta_n)\)。解释为何必须排除 \(n=1\) 的多项式等式。
\end{exercise}
\begin{solution}
对奇数 \(n\)，本原 \(n\) 次根 \(\alpha\) 的负数具有阶 \(2n\)：其 \(n\) 次幂为 \(-1\)，平方具有阶 \(n\)。这个操作为两类本原根建立双射。又因 \(n>1\) 为奇数，互素剩余类由 \(a\) 与 \(-a\) 成对，\(\varphi(n)\) 为偶数，所以 \(\Phi_n(-X)\) 仍首一，得到多项式等式。域方面，不妨选取兼容的本原根 \(\zeta_n=\zeta_{2n}^2\)；此时 \(\zeta_{2n}=-\zeta_n^{(n+1)/2}\)，故两边包含。\(n=1\) 时 \(\Phi_1(-X)=-X-1=-\Phi_2(X)\)，符号不同。
\end{solution}

\begin{exercise}\label{b1:ex:21-3}
证明 \(\mathbb Q(\zeta_8)=\mathbb Q(i,\sqrt2)\)，列出其全部二次子域。
\end{exercise}
\begin{solution}
由 \(\zeta_8=(1+i)/\sqrt2\)，得到左域包含于右域。另一方面 \(i=\zeta_8^2\)，且 \(\sqrt2=\zeta_8+\zeta_8^{-1}\)，给出反向包含。群 \((\mathbb Z/8\mathbb Z)^\times\) 有四个元素，非单位元阶均为 \(2\)，故为 \(V_4\)，恰有三个二阶子群。相应三个二次子域为 \(\mathbb Q(i)\)、\(\mathbb Q(\sqrt2)\)、\(\mathbb Q(\sqrt{-2})\)；它们分别由三种非零平方类生成，彼此不同，已经穷尽。
\end{solution}

\begin{exercise}\label{b1:ex:21-4}
令 \(t=\zeta_5+\zeta_5^{-1}=2\cos(2\pi/5)\)。求 \(t\) 的最小多项式及 \(\cos(2\pi/5)\) 的根式表达。
\end{exercise}
\begin{solution}
将 \(1+\zeta_5+\cdots+\zeta_5^4=0\) 除以 \(\zeta_5^2\)，得 \((t^2-2)+t+1=0\)，即 \(t^2+t-1=0\)。它没有有理根，所以不可约。因为 \(2\pi/5\) 位于第一象限，\(t>0\)，故 \(t=(-1+\sqrt5)/2\)，于是 \(\cos(2\pi/5)=(\sqrt5-1)/4\)。
\end{solution}

\begin{exercise}\label{b1:ex:21-5}
分别取 \(K=\mathbb Q(\sqrt5),\ n=5\) 和 \(K=\mathbb Q(i),\ n=8\)。求 \(K(\mu_n)/K\) 的次数及其 Galois 群，并写出它在 \((\mathbb Z/n\mathbb Z)^\times\) 中对应的子群。
\end{exercise}
\begin{solution}
对第一组数据，\cref{b1:ex:21-4}中的计算给出 \(\sqrt5=2(\zeta_5+\zeta_5^{-1})+1\)，故 \(K\subseteq\mathbb Q(\zeta_5)\)，而 \(K(\mu_5)=\mathbb Q(\zeta_5)\)。两域在 \(\mathbb Q\) 上的次数分别为 \(2\)、\(4\)，所以相对次数为 \(2\)。复共轭固定 \(\sqrt5\)，送 \(\zeta_5\) 到 \(\zeta_5^{-1}\)；它与恒等自同构已穷尽相对 Galois 群。因此\cref{b1:prop:cyclotomic-general-base}的嵌入将这个 \(C_2\) 识别为 \((\mathbb Z/5\mathbb Z)^\times\) 中的 \(\{\overline1,\overline4\}\)。

对第二组数据，\(i=\zeta_8^2\)，故 \(K(\mu_8)=\mathbb Q(\zeta_8)\)，相对次数也是 \(4/2=2\)。自同构 \(\zeta_8\mapsto\zeta_8^a\) 固定 \(i\) 当且仅当 \(i^a=i\)，即 \(a\equiv1\pmod4\)。相应子群为 \(\{\overline1,\overline5\}\)，也同构于 \(C_2\)。这两个例子中的嵌入均非满射：扩大基域后，必须保留新基域元素满足的条件，允许的单位根置换随之减少。
\end{solution}

\begin{exercise}\label{b1:ex:21-6}
用三次求根的配对规则求 \(y^3-3y-2=0\) 的全部根。若独立选择 \(u^3=v^3=1\)，为何不一定得到解？
\end{exercise}
\begin{solution}
此时 \(\Delta=0\)，\(U=V=1\)，但还必须有 \(uv=1\)。选 \(u=v=1\)，配对所得根为 \(2\)、\(\zeta_3+\zeta_3^2=-1\)、\(-1\)，与 \((y-2)(y+1)^2\) 一致。若错误地选 \(u=1,v=\zeta_3\)，则 \(uv=\zeta_3\neq1\)，展开式留下 \((3uv-3)(u+v)\neq0\)，所得和不是原方程的根。
\end{solution}

\begin{exercise}\label{b1:ex:21-7}
对 \(f(X)=X^4+4X-1\)，找出一个二次扩张 \(E/\mathbb Q\)，使 \(f\) 在 \(E[X]\) 中分解为两个二次因子，并写出分解。再说明 \(E/\mathbb Q\) 的非平凡自同构如何作用于这两个因子。
\end{exercise}
\begin{solution}
此处一般四次式的系数为 \(b=c=0\)、\(d=4\)、\(e=-1\)。因三次项已经为零，不需要平移；写成 \(X^4+pX^2+qX+r\) 时，系数就是 \(p=0,q=4,r=-1\)。故辅助多项式为
\[
R(T)=T^3+4T-16=(T-2)(T^2+2T+8).
\]
取辅助根 \(T=2\)，对应的配方参数为 \(m=1\)。由 \(s^2=2m-p=2\)，可取 \(s=\sqrt2\)，而 \(q/(2s)=\sqrt2\)。因此
\[
X^4+4X-1
=\bigl(X^2-\sqrt2X+1+\sqrt2\bigr)
 \bigl(X^2+\sqrt2X+1-\sqrt2\bigr).
\]
这里可以取 \(E=\mathbb Q(\sqrt2)\)。其非平凡自同构把 \(\sqrt2\) 变为 \(-\sqrt2\)，因而交换这两个二次因子。
\end{solution}

\begin{exercise}\label{b1:ex:21-8}
求 \(X^4-2\) 在 \(\mathbb Q\) 上的分裂域及其群，说明它根式可解而 Galois 群非交换。
\end{exercise}
\begin{solution}
令 \(\alpha=\sqrt[4]{2}>0\)，分裂域为 \(L=\mathbb Q(\alpha,i)\)。Eisenstein 判别给出 \(\bigl[\mathbb Q(\alpha):\mathbb Q\bigr]=4\)；该域为实域，加入 \(i\) 再扩张两倍，故 \([L:\mathbb Q]=8\)。在 \(\mathbb Q(i)\) 上，\(\alpha\) 的最小多项式仍为四次，故映射 \(r(\alpha)=i\alpha,r(i)=i\) 定义四阶自同构。复共轭 \(s\) 满足 \(s(\alpha)=\alpha,s(i)=-i\)，并有 \(s^2=1\)、\(srs=r^{-1}\)。八个映射 \(r^j,r^js\) 两两不同，所以群为二面体群 \(D_8\)，本书的下标表示其阶为 \(8\)。塔 \(\mathbb Q\subseteq\mathbb Q(\alpha)\subseteq L\) 直接给出根式表达，而 \(sr\neq rs\)。
\end{solution}

\begin{exercise}\label{b1:ex:21-9}
设特征零域 \(K\) 上的两个多项式 \(f,g\) 都根式可解。证明 \(fg\) 也根式可解，并用 Galois 群给出另一说明。
\end{exercise}
\begin{solution}
先取包含 \(f\) 全部根的根式塔，再把一条为 \(g\) 选定的根式塔逐步与当前末端取合成域。每一步新生成元的某次幂仍在前一层中，所以连接后仍是根式塔，包含两多项式的全部根。群论方面，设分裂域为 \(L_f,L_g\)，则合成域 \(L_fL_g\) 的 Galois 群通过两次限制嵌入 \(\operatorname{Gal}(L_f/K)\times\operatorname{Gal}(L_g/K)\)。这里限制映射的核是平凡群：若一个自同构在 \(L_f\)、\(L_g\) 上都恒等，它便固定二者生成的整个合成域。两个因子可解，有限直积及其子群可解，故\cref{b1:thm:radical-solvability}再次给出结论。
\end{solution}

\begin{exercise}\label{b1:ex:21-10}
求 \(X^5-2\) 在 \(\mathbb Q\) 上的 Galois 群，并与一般五次方程比较。
\end{exercise}
\begin{solution}
令 \(\alpha=\sqrt[5]{2}\)、\(F=\mathbb Q(\zeta_5)\)。Eisenstein 判别给出 \(\bigl[\mathbb Q(\alpha):\mathbb Q\bigr]=5\)，而 \([F:\mathbb Q]=4\)。若 \(\alpha\in F\)，则 \(\mathbb Q(\alpha)\subseteq F\)，塔公式会给出 \(5=[\mathbb Q(\alpha):\mathbb Q]\mid[F:\mathbb Q]=4\)，矛盾。因此 \(\alpha\notin F\)。因 \(F\) 含 \(\mu_5\)，\cref{b1:prop:single-radical-cyclic}说明 \(\bigl[F(\alpha):F\bigr]\) 为 \(1\) 或 \(5\)，只能为 \(5\)。分裂域 \(L=F(\alpha)\) 因而在 \(\mathbb Q\) 上次数为 \(20\)。

定义 \(\sigma(\alpha)=\zeta_5\alpha\)、\(\sigma(\zeta_5)=\zeta_5\)，得到五阶自同构。\(F\) 上的自同构 \(\tau_0\) 由 \(\tau_0(\zeta_5)=\zeta_5^2\) 给定。因 \(X^5-2\) 在 \(F[X]\) 中不可约，映射
\[
F[X]/(X^5-2)\xrightarrow{\ \sim\ }L,\qquad \overline X\longmapsto\alpha
\]
是域同构。通过这一同构，将系数上的 \(\tau_0\) 延拓为 \(L\) 的自同构
\[
\tau\!\left(\sum_{j=0}^4c_j\alpha^j\right)
=\sum_{j=0}^4\tau_0(c_j)\alpha^j\qquad(c_j\in F).
\]
基表示的唯一性保证映射有定义；\(\tau_0\) 固定 \(2\)，所以关系 \(\alpha^5=2\) 得到保持，\(\tau\) 也保持乘法，并且固定 \(\alpha\)。它的阶为 \(4\)。计算得 \(\tau\sigma\tau^{-1}=\sigma^2\)，二者生成全部 \(20\) 个自同构，所以群为非平凡半直积 \(C_5\rtimes C_4\)。它有循环正规子群和循环商，因而可解。这一具体五次方程可解，并不改变独立系数的一般五次方程具有不可解群 \(S_5\) 的结论。
\end{solution}

\begin{exercise}\label{b1:ex:21-11}
为 \(\sqrt2+\sqrt3\) 与 \(\sqrt{1+\sqrt2}\) 分别写出实二次扩张塔，并求它们在 \(\mathbb Q\) 上的次数。
\end{exercise}
\begin{solution}
第一条塔为 \(\mathbb Q\subseteq\mathbb Q(\sqrt2)\subseteq\mathbb Q(\sqrt2,\sqrt3)\)，两个平方类独立，末端次数为 \(4\)。若 \(\alpha=\sqrt3+\sqrt2\)，则 \(\alpha^{-1}=\sqrt3-\sqrt2\)；取和差恢复 \(\sqrt2,\sqrt3\)，所以该和生成末端，次数为 \(4\)。

第二条塔为 \(\mathbb Q\subseteq\mathbb Q(\sqrt2)\subseteq\mathbb Q(\sqrt{1+\sqrt2})\)。\(1+\sqrt2\) 不是 \(\mathbb Q(\sqrt2)\) 中的平方，因为其范数为 \(-1\)，而平方的范数是一个有理数的平方，不可能为负。故第二步也是二次，且 \(\sqrt2=\alpha^2-1\) 属于 \(\mathbb Q(\alpha)\)，其中 \(\alpha=\sqrt{1+\sqrt2}\)。于是 \(\alpha\) 的次数为 \(4\)。两数均由\cref{b1:thm:constructible-quadratic-tower}可作图，因为已经写出了包含它们的实二次塔。这与\cref{b1:prop:constructible-splitting-two-group}后的四次反例形成对照：次数同为 \(4\)，是否可作图仍取决于能否组成这样的塔。
\end{solution}

\begin{exercise}\label{b1:ex:21-12}
已给定角 \(\theta\)，令 \(K=\mathbb Q(\cos\theta,\sin\theta)\)，即其单位圆终边点坐标生成的实域。证明这个角能够尺规三等分，当且仅当 \(X^3-3X-2\cos\theta\) 在 \(K[X]\) 中可约。
\end{exercise}
\begin{solution}
该三次多项式的三个根为 \(2\cos\bigl((\theta+2j\pi)/3\bigr)\)、\(j=0,1,2\)。若不可约，每个根在 \(K\) 上次数为 \(3\)，由相对于 \(K\) 的\cref{b1:thm:constructible-quadratic-tower}，任何一个都不可作出，因此不能三等分。若可约，三次多项式有一个根在 \(K\) 中，除去该线性因子后其余两个根仍是上面列出的实余弦值，所以剩余二次因子的判别式非负，两根在至多二次的实扩张中。于是全部三个余弦可作出，相应正弦由开平方得到，能从所得有限个终边中选出位于给定角内的三等分线。重根时同样适用，不需要三根互异。
\end{solution}

\begin{exercise}\label{b1:ex:21-13}
判断正 \(9,21,34,60,85\) 边形能否尺规作图，并逐一说明理由。
\end{exercise}
\begin{solution}
\(9=3^2\) 含奇素数的重复因子，\(\varphi(9)=6\)，不可作图。\(21=3\cdot7\)，其中 \(7-1=6\) 不是 \(2\) 的幂，\(\varphi(21)=12\)，也不可作图。其余分解为 \(34=2\cdot17\)、\(60=2^2\cdot3\cdot5\)、\(85=5\cdot17\)，奇素因子均为不同的 Fermat 素数。相应 Euler 函数值为 \(16,16,64\)，由\cref{b1:thm:constructible-polygon}都可作图。
\end{solution}

\begin{exercise}\label{b1:ex:21-14}
设 \(n\geq3\)。证明 \(\Bigl[\mathbb Q\bigl(2\cos(2\pi/n)\bigr):\mathbb Q\Bigr]=\varphi(n)/2\)，并据此判断 \(\cos(2\pi/7)\) 能否作图。
\end{exercise}
\begin{solution}
记 \(t=\zeta_n+\zeta_n^{-1}\)。\(\zeta_n\) 满足 \(X^2-tX+1=0\)，且非实，而 \(\mathbb Q(t)\subseteq\mathbb R\)，所以 \(\bigl[\mathbb Q(\zeta_n):\mathbb Q(t)\bigr]=2\)。分圆域次数为 \(\varphi(n)\)，由塔公式得所需等式。对 \(n=7\)，\(t\) 的次数为 \(3\)，乘以有理常数 \(1/2\) 不改变生成域，因此 \(\cos(2\pi/7)\) 不可作图。
\end{solution}

\begin{exercise}\label{b1:ex:21-15}
从单位线段出发，能否作出一条长度等于单位圆周长的线段？这与化圆为方所用的障碍是否相同？
\end{exercise}
\begin{solution}
单位圆周长为 \(2\pi\)。若能作出该线段，规矩数域对除以 \(2\) 封闭，所以 \(\pi\) 也可作图，因而代数。这与由 Lindemann--Weierstrass 定理推出的\cref{b1:cor:pi-transcendental}矛盾，故不能。化圆为方要求 \(\sqrt\pi\) 可作图，平方后同样会迫使 \(\pi\) 代数；两者都依赖圆周率的超越性，不能仅由某个有限扩张次数的计算代替。
\end{solution}


\begin{exercise}\label{b1:ex:21-reduction-quartic}
令 \(g(X)=X^4-6X-2\)，\(L\) 为它在 \(\mathbb Q\) 上的分裂域。
\begin{enumerate}
\item 证明 \(g\) 不可约，核验分解
\[
\begin{aligned}
g(X)&\equiv(X-1)(X^3+X^2+X+2)\pmod7,\\
g(X)&\equiv(X-6)(X-8)(X^2-3X+12)\pmod{17},
\end{aligned}
\]
并据此求 \(\operatorname{Gal}(L/\mathbb Q)\)。
\item 证明 \(g\) 在 \((1,2)\) 中有实根，判断这个根能否用根式表达、能否尺规作出。
\item 已知 \(\operatorname{disc}(g)=-37040\)，求 \(L\) 中的二次子域，并证明它是唯一的。
\end{enumerate}
\end{exercise}
\begin{solution}
1. 多项式 \(g\) 对素数 \(2\) 满足 Eisenstein 条件，因而在 \(\mathbb Q[X]\) 中不可约。展开乘积可核验题中的两处分解。模 \(7\) 的三次因子在 \(0,1,\ldots,6\) 处的值依次为 \(2,5,2,6,2,3,1\)，没有根，故不可约。模 \(17\) 的二次因子的判别式为 \(9-48\equiv12\)，而非零平方剩余为 \(1,2,4,8,9,13,15,16\)，所以它不可约；两个一次根 \(6,8\) 互异，也不是二次因子的根。两次约化均无重根。

由\cref{b1:lem:reduction-cycle-type}，群 \(G=\operatorname{Gal}(L/\mathbb Q)\leq S_4\) 含三轮换和换位；不可约性又使它在四根上传递。因此 \(4\bigm\vert|G|\) 且 \(3\bigm\vert|G|\)，只能有 \(|G|=12\) 或 \(24\)。若阶为十二，指数为二使 \(G\) 在 \(S_4\) 中正规，因而包含所含换位的全部共轭，也就是所有换位；但换位生成 \(S_4\)，矛盾。故 \(G=S_4\)。只用三轮换不能得到这一结论，因为传递子群 \(A_4\) 也包含三轮换。

2. 因 \(g(1)=-7\)、\(g(2)=2\)，连续性给出一个实根 \(\alpha\in(1,2)\)。群 \(S_4\) 有正规列 \(1\subset V_4\subset A_4\subset S_4\)，各因子为交换群，故由\cref{b1:thm:radical-solvability}，\(\alpha\) 能用根式表达。但 \(S_4\) 的阶为 \(24\)，不是 \(2\) 的幂；\cref{b1:prop:constructible-splitting-two-group}排除了尺规作图。不可约性给出 \([\mathbb Q(\alpha):\mathbb Q]=4\)，所以这个计算再次区分了代数次数的条件与作图条件。

3. 将四根记为 \(a_1,\ldots,a_4\)，置 \(\delta=\prod_{i<j}(a_j-a_i)\)。\cref{b1:prop:discriminant-alternating-group}的证明给出 \(\sigma(\delta)=\operatorname{sgn}(\sigma)\delta\)，故 \(\delta\) 的稳定子群恰为 \(A_4\)。于是
\[
L^{A_4}=\mathbb Q(\delta)
=\mathbb Q(\sqrt{-37040})
=\mathbb Q(\sqrt{-2315}).
\]
这里 \(-37040=16\cdot(-2315)\)，且负有理数不是有理数的平方，故这个子域确为二次。若还有二次子域，Galois 对应会给出 \(S_4\) 的另一个指数为二的子群，也就是某个满同态 \(S_4\longrightarrow C_2\) 的核。所有换位共轭，在该同态下有相同的像；换位生成 \(S_4\)，满射性迫使这个像为 \(C_2\) 的非单位元。因此该同态就是符号同态，其核只能是 \(A_4\)。这证明了唯一性。
\end{solution}
